% Compile with pdfLaTeX. Run twice to resolve internal references.
\documentclass[11pt,a4paper,reqno]{amsart}
\usepackage[T1]{fontenc}
\usepackage{lmodern,mathtools,amssymb,booktabs,microtype}
\usepackage[margin=27mm,top=26mm,bottom=26mm]{geometry}
\usepackage{tikz}
\usepackage[unicode,bookmarksnumbered,bookmarksdepth=2]{hyperref}
\definecolor{ink}{HTML}{214555}
\hypersetup{colorlinks=true,linkcolor=ink,citecolor=ink,urlcolor=ink,
  pdftitle={The polygonal Pólya–Szegő inequality for the first Dirichlet eigenvalue},
  pdfauthor={},pdfsubject={A unified analytic proof for simple planar polygons}}
\newtheorem{theorem}{Theorem}[section]
\newtheorem{lemma}[theorem]{Lemma}
\newtheorem{proposition}[theorem]{Proposition}
\theoremstyle{remark}
\newtheorem*{remark}{Remark}
\numberwithin{equation}{section}
\newcommand{\dd}{\,\mathrm d}
\newcommand{\av}[1]{\overline{#1}}
\newcommand{\davg}[1]{\langle #1\rangle}
\newcommand{\norm}[1]{\left\|#1\right\|_2}
\newcommand{\Cphi}{C_\phi}
\newcommand{\Jcal}{\mathcal J}
\newcommand{\D}{\mathbb D}
\newcommand{\R}{\mathbb R}
\DeclareMathOperator{\Var}{Var}
\allowdisplaybreaks[1]
\setcounter{tocdepth}{1}
\emergencystretch=1.5em
\raggedbottom

\title[The first Dirichlet eigenvalue of a polygon]{The polygonal P\'olya--Szeg\H{o} inequality\\ for the first Dirichlet eigenvalue}
\date{November 22, 2026}
\subjclass[2020]{35P15, 49Q10, 35J05, 30C20}
\keywords{First Dirichlet eigenvalue, polygonal Faber--Krahn inequality, conformal mapping, shape optimization, superlevel sets}
\begin{document}
\begin{abstract}
Among simple planar polygons of prescribed area with at most $N$ sides, the regular
$N$-gon uniquely minimizes the first Dirichlet eigenvalue, up to congruence.
We give an analytic proof that treats every $N\ge3$ by the same argument.
For convex polygons, a partition into mixed boundary value problems reduces
the comparison to a global supporting-line inequality for isosceles triangles.
A bound in terms of the number of convex vertices then reduces any putative
minimizing sequence below the regular value to polygons with one reflex vertex.
Compactness in that class yields a minimizer. At a nonconvex stationary polygon,
the boundary pressure forces the ground state to have a disconnected superlevel
set near its root-mean-square height. For an area-normalized eigenvalue no larger than that of the regular pentagon,
a quantitative estimate on the disk gives connected superlevel sets throughout
a fixed neighborhood of the same height. These two conclusions exclude nonconvex minimizers.
\end{abstract}
\maketitle


\section{Introduction}\label{sec:main}

The polygonal P\'olya--Szeg\H{o} problem asks which polygon of a prescribed
area and side count has the smallest principal frequency. In the triangle and
quadrilateral cases the extremizers are the equilateral triangle and the square
\cite[pp.~158--159]{ps}. For more sides a Steiner symmetrization can increase the
side count. Wang's preprint states the sharp eigenvalue comparison for all
simple polygons with at most $N$ sides when $N\ge5$, together with results on
torsional rigidity and stability \cite[Theorem~1.1]{wang}. Here we focus on the
eigenvalue and give a proof with a single range of comparison angles for all
$N\ge3$.

A \emph{simple polygon} $P$ is the bounded open region enclosed by a polygonal
Jordan curve. Merge adjacent collinear sides and denote the number of remaining
sides by $n(P)$. A vertex is \emph{convex} if its interior angle is less than
$\pi$, and \emph{reflex} if it is greater than $\pi$. Write $p(P)$ and $q(P)$ for
their respective numbers, so that $p(P)+q(P)=n(P)$. For a regular $N$-gon
$R_N$ of arbitrary size, set
\[
 \lambda_1(P)=\inf_{0\ne u\in H_0^1(P)}
       \frac{\int_P |\nabla u|^2}{\int_P |u|^2},\qquad
 J(P)=|P|\lambda_1(P),\qquad r_N=J(R_N),
\]
The functional $J$ is invariant under dilations.

\begin{theorem}\label{thm:main}
Let $N\ge3$ be an integer. Every simple polygon $P$ with $n(P)\le N$ satisfies
\begin{equation}\label{eq:main}
                         J(P)\ge r_N .
\end{equation}
Equality holds precisely when $P$ is similar to $R_N$. In particular, at fixed
area $A>0$ the minimum is $r_N/A$ and the minimizer is unique up to rigid
motions, including reflection. This holds both for exactly $N$ and for at most
$N$ effective sides.
\end{theorem}

Two distinct issues arise in passing from convex to simple polygons. Even
if every minimizer is convex, a minimizing sequence might leave the class of
simple polygons in the limit. We first show that a polygon below $r_N$ would
have exactly one reflex vertex. Cutting at that vertex yields two convex
pieces and recovers a simple competitor in the limit, with control of both area
and side count. We can then use stationarity to prove convexity of the attained
minimizer.

For the convex comparison, Section~\ref{sec:angles} develops Wang's mixed-cell
inequality and balanced partition \cite[Lemmas~2.1 and~2.2]{wang}. The common
opening of the regular cells is $2\pi/N$. We prove that the tangent to their
spectral function is a strict global supporting line at every opening up to
$2\pi/3$. The proof uses a sector near zero, the square spectrum near $\pi/2$,
and the equilateral triangle at $2\pi/3$. In Section~\ref{sec:turns}, Wang's
conformal coefficient identities and shifted Jensen inequality
\cite[Lemmas~3.1--3.3]{wang} give a bound for polygons with at most $N-2$ convex
vertices. The extension of the supporting-line argument, the common integral
weight for the turning measure and the compactness argument of
Section~\ref{sec:compact} follow the working manuscript \cite{unified}.

Sections~\ref{sec:pressure}--\ref{sec:levels} develop the convexity argument
in the manuscript made available by Zhuo Cheng \cite{cheng}. A stationary nonconvex
polygon must have at least five sides. Thus every such minimizer has the regular
pentagon available as a competitor. This gives a single spectral threshold at
which the boundary pressure count and the disk estimates apply simultaneously.
Section~\ref{sec:assembly} finishes the proof and identifies the equality case.
Appendices~\ref{app:acute}--\ref{app:series} contain the estimates for the sector
and the disk. We use the classical Faber--Krahn and Kohler--Jobin inequalities,
the Schwarz--Christoffel formula, and standard rearrangement and spectral
perturbation theory. All terminating decimals in estimates denote rational numbers.

\tableofcontents

\section{The sharp inequality for convex polygons}\label{sec:angles}
\subsection{A mixed problem for each side}\label{sec:cell}

Let $I_\alpha$ be the isosceles triangle of altitude one and apex angle
$\alpha\in(0,\pi)$. For a polygon $Q$ with a designated base, denote by
$\mu_D(Q)$ the lowest eigenvalue with Dirichlet condition on the base and
Neumann condition on the rest of the boundary. Put
\[
 F(\alpha)=|I_\alpha|\mu_D(I_\alpha)=\tan(\alpha/2)\mu_D(I_\alpha).
\]
Every line perpendicular to the base cuts $I_\alpha$ in a segment of length
at most one. The one-dimensional mixed Poincar\'e inequality gives
$\mu_D(I_\alpha)\ge\pi^2/4$. Testing with distance to the base gives
$\mu_D(I_\alpha)\le6$. In particular $F(0+)=0$ and $F(\pi-)=+\infty$.
An affine change of variables to a fixed triangle gives continuity in between.

\begin{lemma}[Mixed-cell comparison, Wang]\label{lem:cell}
Suppose the two angles at the base of a convex polygon $Q$ are $\pi a$ and
$\pi b$, with $C=a+b<1$. For $\alpha=\pi(1-C)$ one has
\[
                   |Q|\mu_D(Q)\ge F(\alpha).
\]
Equality holds precisely for an isosceles triangle on the designated base.
\end{lemma}

\begin{proof}
Map the upper half-plane conformally onto $Q$ with $0$ and $\infty$
corresponding to the ends of the base and the positive real axis to its
interior. Write the other prevertices as $-e^{x_j}$. Their exterior angles
are $\pi c_j$, where $c_j>0$ and $\sum c_j=C$. The substitution
$z=e^{x+iy}$ identifies the half-plane with the strip
$S=\R\times(0,\pi)$. By the Schwarz--Christoffel formula \cite{SC}, the
pullback of area has density proportional to
\[
       e^{\delta x}\prod_j[\cosh(x-x_j)+\cos y]^{-c_j},
       \qquad \delta=a-b .
\]
Normalize this density to a probability density $p$. Then
$(|Q|\mu_D(Q))^{-1}$ is the supremum of $\int_S p u^2$ over nonnegative
functions of unit Dirichlet energy and zero bottom trace. Indeed,
conformal pullback preserves the energy and the zero trace locally on the
open base. Integrating from that trace along vertical lines bounds the
unweighted strip $L^2$ norm by the energy. Conversely, the density has
exponentially decaying tails since $|\delta|<C$, and each corner singularity
has order less than two. Thus $p\in L^r(S)$ for some $r>1$, and
H\"older's inequality and the strip Sobolev embedding give the weighted
integrability needed for the inverse change of variables. The mixed problem
on the bounded polygon has a ground state.

Represent each denominator by its gamma integral, put $v_j=vt_j$ with
$\sum_j t_j=1$, and integrate out $v$. The remaining integral has simplex
weight proportional to $\prod_j t_j^{c_j-1}$ and denominator
$[\sum_jt_j\cosh(x-x_j)+\cos y]^C$. Write the horizontal sum as
$K\cosh(x-\bar x)$, where
\[
 K^2=\left(\sum_jt_je^{x_j}\right)\left(\sum_jt_je^{-x_j}\right)\ge1 .
\]
After normalizing the components, $p$ is a probability mixture of horizontal
translates of
\[
 p_{\delta,\kappa}\ \propto\ e^{\delta x}(\cosh x+\kappa\cos y)^{-C},
 \qquad \kappa=K^{-1}.
\]
If two $x_j$ are distinct, then $\kappa<1$ throughout the open simplex.
Horizontal translations preserve the trial space, so we can compare each
component separately.

Rearrange a trial function first in increasing order on each vertical slice,
then in symmetric decreasing order on each horizontal slice. Call the result
$v$. For the vertical step, reflect evenly about $y=\pi$, extend by zero, and
rearrange about $\pi$. Rearrangement decreases the energy in the rearranged
direction. In the transverse direction apply its $L^2$ contraction property to
difference quotients of adjacent slices. The bottom trace stays zero, and the
horizontal step preserves vertical order. Hence $v$ is nondecreasing in $y$
and nonincreasing in $|x|$. The Hardy--Littlewood inequality bounds the original moment by its vertically
rearranged value since $\partial_y\log p_{\delta,\kappa}
=C\kappa\sin y/(\cosh x+\kappa\cos y)\ge0$.

For the horizontal step, expand in nonnegative powers of
$(1-\kappa\cos y)/(1+\cosh x)$. The density is a mixture
$\sum_{k\ge0}w_{k,\delta,\kappa}\rho_{C+k,\delta}(x)\eta_{k,\kappa}(y)$,
with probability weights and one-dimensional probability densities
\[
 \rho_{d,\delta}\propto e^{\delta x}(1+\cosh x)^{-d},\qquad
 \eta_{k,\kappa}\propto(1-\kappa\cos y)^k .
\]
The Hardy--Littlewood inequality bounds each horizontal moment by the
corresponding moment of $v^2$ against the symmetric decreasing rearrangement
$\rho_{d,\delta}^*$.

We compare this rearranged density with $\rho_{d,0}$. Let $[m-h,m+h]$ be a
superlevel interval of the original density, and let $L,R$ be the absolute
values of its logarithmic derivative at the two endpoints. As the logarithm
of the level increases, the endpoint velocities are $1/L$ and $-1/R$.
The absolute logarithmic derivative of the rearranged density at $h$ is
therefore $2LR/(L+R)$. Write $r=\tanh(h/2)$ and $z=\tanh(m/2)$. The
original logarithmic derivative is $\delta-d\tanh(x/2)$, so
\[
 \frac{2LR}{L+R}\le\frac{L+R}{2}
 =dr\frac{1-z^2}{1-z^2r^2}\le dr
 =-\partial_h\log\rho_{d,0}(h).
\]
Thus $\rho_{d,\delta}^*/\rho_{d,0}$ increases with $h$. Since both
densities have mass one, $\rho_{d,0}$ gives at least as much mass to every
centered interval. If $\delta\ne0$, equality of the endpoint values forces
$m\ne0$, and the comparison is strict on every centered interval of positive
finite length. Layer integration against $v^2$ gives strict gain for a
nonzero trial function.

Set $a_k(\kappa)=\int_S\rho_{C+k,0}\eta_{k,\kappa}v^2$. As $k$ increases,
ratios of successive densities concentrate the horizontal mass at $x=0$ and
move the vertical mass upward. Hence $a_k(\kappa)$ is nondecreasing in $k$.
The horizontal normalizing integral is $2^d\mathrm B(d+\delta,d-\delta)$.
Successive values of $w_{k,\delta,\kappa}/w_{k,0,\kappa}$ have ratio
$1-\delta^2/(C+k)^2$. Replacing the weights by $w_{k,0,\kappa}$ therefore
increases the average. At fixed $k$, increasing $\kappa$ gives a vertical
likelihood ratio increasing in $X=-\cos y$. Thus $a_k(\kappa)\le a_k(1)$.
For $k=1$ the gain is strict when $\kappa<1$, since every nonzero vertical
slice of $v$ is nonconstant and has zero bottom trace. Finally, set
$Z_k(\kappa)=\int_0^\pi(1+\kappa X)^k\dd y$. Then
\[
 Z_{k+1}(\kappa)/Z_k(\kappa)=1+\kappa\mathbb E_{k,\kappa}X .
\]
where the expectation is under $\eta_{k,\kappa}$. Symmetry gives
$\mathbb E_{k,\kappa}X\ge0$, and the likelihood-ratio ordering makes it
nondecreasing in $\kappa$. Consequently $w_{k,0,1}/w_{k,0,\kappa}$ is
nondecreasing in $k$. All these comparisons of averages follow from
nonnegative covariance of functions monotone in the same direction, with
truncation if needed.

Each trial moment is now bounded by the optimum for $p_{0,1}$, the density
of $I_\alpha$. Restoring unit energy after rearrangement can only increase
the moment. Apply the strict comparisons to the maximizer on $Q$ and
integrate over the simplex. Equality requires $\delta=0$ and exactly one
prevertex off the base. These conditions give the stated isosceles triangle.
\end{proof}

\subsection{Balancing the cell areas}\label{sec:balance}

Give each side of a convex polygon $P$ a positive weight $w_i$. Let $Q_i$
consist of the points for which $w_i$ times the distance to the supporting
line of side $i$ is least. The cells are convex and their interiors are disjoint. At a relative interior
point of side $i$ its weighted distance is zero while all the others are
positive. Thus $Q_i$ occupies a neighborhood of that point in $P$ and has
positive area. At vertex $i$, let $\theta_i$ be the angle
between sides $i-1$ and $i$, and let $t_i$ be the angle between the separating
ray and side $i-1$. If $\pi-\alpha_i$ is the sum of the base angles of $Q_i$,
then
\[
 \frac{\sin t_i}{\sin(\theta_i-t_i)}=\frac{w_i}{w_{i-1}},
 \qquad \alpha_i=\pi-\theta_i+t_i-t_{i+1},\qquad \sum_i\alpha_i=2\pi .
\]
The angles vary continuously with the weights. The areas are continuous as
well, because every tie between two weighted distance functions has area zero.

Extend $F$ by zero to nonpositive openings. For $\epsilon>0$, take
$w_i=(\omega_i+\epsilon)^{-1}$ on the closed probability simplex and put
$G_i=F(\alpha_i)/\sum_j F(\alpha_j)-|Q_i|/|P|$. The sum of the openings is
positive, so the denominator is positive. The continuous simplex map
\[
 \omega_i\longmapsto\frac{\omega_i+[G_i]_+}{1+\sum_j[G_j]_+}
\]
has a Brouwer fixed point. There $[G_i]_+=\omega_i\sum_j[G_j]_+$ and
$\sum_j[G_j]_+\le1$.

Pass to a convergent subsequence as $\epsilon\downarrow0$, and let $Z$ be the
indices whose limiting coordinates vanish. If $Z$ is nonempty, the weighted
distances for these indices tend to infinity at every interior point of $P$.
Comparison with any index outside $Z$ shows that the total area of their
cells tends to zero. The fixed-point identity then forces their total target
area fraction to tend to zero. For $i\notin Z$ the weight $w_i$ is bounded
above and $w_{i-1}$ bounded below. Thus $t_i$ stays away from $\theta_i$,
$\alpha_i$ stays away from $\pi$, and the sum of the $F(\alpha_i)$ over the
complement of $Z$ is bounded. On any maximal cyclic block $k,\ldots,l$ in
$Z$ we have $t_k\to\theta_k$ and $t_{l+1}\to0$. Its openings sum in the
limit to $\sum_{i=k}^l(\pi-\theta_i)+\theta_k\ge\pi$. At least one opening
in the block is bounded below by a fixed positive number. The bounds on
$F$ then preclude a vanishing target fraction. This proves that the limiting
point lies in the interior. If one limiting $G_i$ were positive, the
fixed-point identity would make every $G_i$ positive, contrary to $\sum_iG_i=0$.
Thus all $G_i$ vanish, and Lemma~\ref{lem:cell} gives
\begin{equation}\label{eq:cells}
 \alpha_i>0,\quad |Q_i|=|P|\frac{F(\alpha_i)}{\sum_j F(\alpha_j)},\quad
 |Q_i|\mu_D(Q_i)\ge F(\alpha_i).
\end{equation}
Restrict a Dirichlet trial function on $P$ to the cells, apply the mixed
Rayleigh inequality on each, and sum. We obtain
\begin{equation}\label{eq:sumF}
                    J(P)\ge\sum_i F(\alpha_i),\qquad r_N=N F(2\pi/N).
\end{equation}
For the identity on the right, the regular ground state is even across each
radial symmetry line. Conversely the isosceles mixed ground states match
across these lines by reflection.

\begin{proposition}[Supporting lines]\label{prop:support}
The function $F$ satisfies $F''>0$ on $(0,2\pi/3]$. For $0<a\le2\pi/3$ and
$0<\alpha<\pi$ with $\alpha\ne a$,
\begin{equation}\label{eq:fulltangent}
              F(\alpha)>F(a)+F'(a)(\alpha-a).
\end{equation}
\end{proposition}

Appendix~\ref{app:acute} proves positivity of the curvature for
$0<\alpha\le2\pi/5$ by a conformal change to a sector.
Section~\ref{sec:square} proves $F''>1/10$ on
$[2\arctan(1/\sqrt2),\pi/2]$. These intervals overlap because
$\cos(2\pi/5)=(\sqrt5-1)/4<1/3$. We establish the remaining curvature bound
and the global support of the tangent at $2\pi/3$ below.

\subsection{The equilateral ground state and its spectral gap}\label{sec:equispec}
Cut the isosceles cell along its symmetry axis, across which the ground state
is even. Use the coordinates
\[
 t=\tan(\alpha/2),\quad z=t^{-2},\quad
 D_0=\{0<y<x<1\},\quad \dd\mu=2\dd x\dd y .
\]
The physical point is $(x,ty)$. On functions with zero trace on $x=1$, let
$A$ and $B$ be the quadratic forms $\int u_x^2\dd\mu$ and $\int u_y^2\dd\mu$.
The other two sides carry natural conditions. Denote by $G(z)$ the lowest
eigenvalue of $Q_z=A+zB$ in $L^2(\dd\mu)$. Then $F(\alpha)=tG(z)$. The forms
have a common domain and depend analytically on $z$. By simplicity of the
ground state, $G$ is analytic for $z>0$ \cite{kato}. Horizontal slices have
length at most one, so $A\ge\ell\|\cdot\|_2^2$ with $\ell=\pi^2/4$.

Put $c=1/3$, $H=4\pi^2/9$ and $\Lambda=4H$, so $\ell=9H/16$. At $z=c$ the
physical half-cell reflects to an equilateral triangle of inradius one
(Figure~\ref{fig:equilateral}). The distances $d_i$ to its three sides sum
to three. The function
$\sum_i\sin(2\pi d_i/3)=4\prod_i\sin(\pi d_i/3)$ is positive inside and
vanishes on the boundary. Each summand satisfies the eigenvalue equation
with eigenvalue $H$. Thus $G(c)=H$. In our coordinates the eigenfunction
and its unit normalization are
\[
 E(x,y)=\sin\!\left(\tfrac{2\pi}{3}(1-x)\right)
    +2\sin\!\left(\tfrac{2\pi}{3}+\tfrac{\pi x}{3}\right)\cos(\pi y),
 \qquad e=E/\sqrt{3/2}.
\]
Integrating over $0<y<x<1$ gives
\begin{equation}\label{eq:e-data}
 \|E\|_2^2=\tfrac32,\quad \|E_y\|_2^2=\pi^2-\tfrac{243}{32},\quad
 \beta=B(e,e)=\tfrac{2\pi^2}{3}-\tfrac{81}{16},\quad
 J_y=\int y e_y\dd\mu=\frac{\sqrt6(8\pi-27\sqrt3)}{16\pi^2}.
\end{equation}
For the last two integrals use $E_y=-2\pi\sin(2\pi/3+\pi x/3)\sin(\pi y)$ and
\[
 \int_0^x\sin^2(\pi y)\dd y=\frac x2-\frac{\sin(2\pi x)}{4\pi},\qquad
 \int_0^x y\sin(\pi y)\dd y=-\frac{x\cos(\pi x)}\pi+\frac{\sin(\pi x)}{\pi^2}.
\]
The bounds $3.141<\pi<3.142$ now give $4.38<H<4.39$, $1.50<\beta<1.52$
and $|J_y|>0.33$.

\begin{figure}[htbp]
\centering
\begin{tikzpicture}[x=0.73cm,y=0.73cm,line join=round,every node/.style={font=\small}]
 \fill[ink!19] (0,0)--(1,0)--(1,1.73205)--cycle;
 \draw[thin,black!55] (0,0)--(-2,0) (0,0)--(1,1.73205) (0,0)--(1,-1.73205);
 \draw[densely dashed,ink] (0,0)--(1,0) (0,0)--(-0.5,0.866025) (0,0)--(-0.5,-0.866025);
 \draw[thick] (-2,0)--(1,1.73205)--(1,-1.73205)--cycle;
 \draw[very thick,ink] (1,0)--(1,1.73205);
 \fill (0,0) circle (1.4pt);
\end{tikzpicture}
\caption{Six copies of the shaded half-cell tile the equilateral triangle.
The heavy side carries the Dirichlet condition. The two sides through the
center carry Neumann conditions.}\label{fig:equilateral}
\end{figure}

On the $L^2$-orthogonal complement $e^\perp$ we claim that
$Q_c\ge\Lambda\|\cdot\|_2^2$. Reflect a mixed eigenfunction evenly in the
two medians to fill the equilateral triangle, then oddly across its three
sides to tile the plane. The reflections preserve the weak equation, traces
match, and isolated vertices are removable for the finite-energy extension.
The extension is invariant under rotations by $120^\circ$ about the center
and about each vertex. The first invariance comes from the two median
symmetries, the second from two side reflections. Composing a vertex
rotation with the inverse center rotation gives translations by differences
of vertices. For vertices $(-2,0),(1,\sqrt3),(1,-\sqrt3)$ the reciprocal
lattice is generated by $(2\pi/3)n_1,(2\pi/3)n_2$, where
\[
 n_1=(1,0),\quad n_2=(-1/2,\sqrt3/2),\quad n_3=-n_1-n_2 .
\]
The squared frequencies divided by $H$ are $h^2-hk+k^2$, $h,k\in\mathbb Z$.
The positive values below four are one and three. The value three corresponds
to $(2\pi/3)(n_i-n_j)$, a wave vector parallel to the remaining side.
Reflection in that side fixes the exponential mode, so odd symmetry forces
its coefficient to vanish. The constant mode also vanishes. The value one
is the simple first eigenvalue, which proves the claimed gap.

\subsection{Positive curvature from the square to the equilateral triangle}
Let $c\le z\le1$, write $u$ for the unit positive ground state of $Q_z$, and
set $\delta=z-c$, $R=(1+z)z^2$, $K=Q_z(u',u')-G\|u'\|_2^2$. Here primes on
$u,G$ denote $z$-derivatives. Differentiating the weak eigenvalue equation
gives $G'=B(u,u)$ and $G''=-2K$. The infimum of the affine Rayleigh
quotients is concave. Using the square ground state
$2\cos(\pi x/2)\cos(\pi y/2)$ at $z=1$ gives
\[
    0\le G'\le\beta,\qquad H\le G\le G(1)=9H/8.
\]
With $k=(1+t^2)/2$, differentiating $tG(t^{-2})$ twice with respect to
$\alpha$ yields
\begin{equation}\label{eq:angcurv}
 \frac{F''(\alpha)}{kt}=G+(z^2-z)G'-4RK .
\end{equation}

Equip $e^\perp$ with the inner product $Q_c-G(\cdot,\cdot)_2$. In this
space let $P,M$ represent the form $B$ and the mass inner product, and let
$f$ represent the functional $B(e,\cdot)$. The gap and the horizontal
Poincar\'e bound give
\[
  0\le P\le3\frac{\Lambda-\ell}{\Lambda-G}I
        \le\frac{165}{46}I<\kappa I,\quad
  0\le M\le\frac I{\Lambda-G},\qquad \kappa=18/5 .
\]
Indeed $B=3(Q_c-A)$, and $(s-\ell)/(s-G)$ decreases for $s\ge\Lambda$.
Cauchy--Schwarz for $B$ gives $f=P^{1/2}\xi$ for some $\|\xi\|^2\le\beta$.
Write $C=(I+\delta P)^{-1}$. The projected equation has the solution
$u=\gamma(e-y)$, where $y=\delta Cf$ and $0<\gamma\le1$.
Differentiate $(Q_z-G)(y,v)=\delta B(e,v)$ on the fixed space $e^\perp$.
This gives
\[
 y'=C^2f+\delta G'CMC f,\qquad K=\gamma^2(Q_z-G)(y',y').
\]
Here subtraction of $G$ denotes subtraction of $G$ times the mass form.
The terms involving $\gamma'$ vanish by the eigenvalue equation. Taking the
$Q_z-G$ norm, using $0<C\le I$, and applying the spectral theorem to $P$
gives
\begin{equation}\label{eq:rbound}
 \sqrt{RK}\le\sqrt\beta
 \left(\sqrt{M_3}+\frac{\delta\beta}{\Lambda-G}\sqrt{M_2}\right),\qquad
 M_j=\max_{0\le s\le\kappa}\frac{R s}{(1+\delta s)^j}.
\end{equation}

We have $M_3\le8/15$ and $M_2\le4/5$ throughout the interval. Indeed, for
$j=2,3$ the maximizing $s$ is $\kappa$ if $\delta\le1/[(j-1)\kappa]$.
Then $z/(1+\kappa\delta)\le1/3$ and
$(1+z)/(1+\kappa\delta)\le4/3$. These inequalities give
$M_3\le4\kappa/27$ and $M_2\le2\kappa/9$. Beyond the respective transition
the maximizing $s$ is $1/[(j-1)\delta]$, giving $4R/(27\delta)$ and
$R/(4\delta)$, respectively. The function $R/\delta$ is convex for
$\delta>0$, as one sees by expanding $(1+c+\delta)(c+\delta)^2/\delta$.
Thus the two endpoints suffice. The transition value was just estimated,
and at $z=1$ the values are $4/9$ and $3/4$.

It follows that $\delta\beta/(\Lambda-G)<0.081$ and
$4RK<4(1.52)(0.731+0.081\cdot.895)^2<3.93$. The other terms in
\eqref{eq:angcurv} satisfy $G+(z^2-z)G'\ge H-\beta/4>4$.
Hence $F''>0$ for $\pi/2\le\alpha\le2\pi/3$.

\subsection{The overlapping interval at the square}\label{sec:square}
Equations \eqref{eq:angcurv} and \eqref{eq:rbound} also apply with the
square as reference, as in \cite[Appendix~S1]{wang}. Set $c=1$,
$H_s=\pi^2/2$ and $\delta=z-1$ for $1\le z\le2$. The reference state
$e_s=2\cos(\pi x/2)\cos(\pi y/2)$ has
$\beta_s=B(e_s,e_s)=H_s/2-2<0.468$, by direct integration. Even reflection
in the diagonal and the coordinate axes expresses the eigenvalues as
multiples of sums of two squared odd frequencies. The first two distinct
sums are two and ten. Hence $Q_1\ge5H_s\|\cdot\|_2^2$ on $e_s^\perp$.
Concavity gives $0\le G'\le\beta_s$ and $H_s\le G\le H_s+\delta\beta_s$.
Since $B=Q_1-A$, the bounds on the projected operators become
\[
 \frac{5H_s-\ell}{5H_s-G}\le\frac{(9/2)H_s}{4H_s-\beta_s}<1.16=\kappa_s,
 \qquad \frac{\delta\beta_s}{5H_s-G}<0.025 .
\]
For $0\le x\le\kappa_s$, $z/(1+\delta x)\le1$ if $x\ge1$. If $x\le1$ this
ratio increases with $z$ and is at most $2/(1+x)$. Use
$R=(1+1/z)z^3=(1+z)z^2$ and maximize $16x/(1+x)^3$ and $4x/(1+x)^2$ on
$[0,1]$. This gives $M_3\le\max(2\kappa_s,64/27)=64/27$ and
$M_2\le3\kappa_s$. The parentheses in \eqref{eq:rbound} are bounded by
$1.54+0.025(1.866)<1.587$. Since $kt\ge3/(4\sqrt2)>0.53$, we obtain
\[
 F''(\alpha)>0.53\,[4.9298-4(0.468)(1.587)^2]>1/10
 \quad\bigl(2\arctan(1/\sqrt2)\le\alpha\le\pi/2\bigr).
\]
This interval overlaps the one in Appendix~\ref{app:acute}.

\subsection{Global support to the right of the equilateral opening}
Put $a_0=2\pi/3$. From $F=tG$ and \eqref{eq:e-data},
\[
             F(a_0)=\sqrt3 H,\qquad
             F'(a_0)=2(H-2c\beta)=27/4 .
\]
We first establish, for $v\in e^\perp$, the coupling estimate
\begin{equation}\label{eq:stream}
              |B(e,v)|^2\le C_* Q_c(v),\qquad C_*<2.6 .
\end{equation}
The vector field $(0,e_y)$ pairs with $\nabla v$ by ordinary integration to
give $B(e,v)$. Use the vector-field inner product
$\int(X_1Y_1+3X_2Y_2)\dd\mu$. Put $\psi=y(x-y)$. Both $(e_x,ce_y)$ and
$(\psi_y,-\psi_x)$ have zero ordinary pairing with $\nabla v$. For the
first this follows from $v\perp e$, and for the second from integration by
parts, $\psi=0$ on $y=0,y=x$, and $v=0$ on $x=1$. The two fields are
orthogonal in the vector-field inner product as well. Their squared norms
are $H,2/3$, and their inner products with $(0,e_y)$ are $\beta,-3J_y$.
Subtract both orthogonal projections and use Cauchy--Schwarz against the
dual gradient norm $Q_c(v)^{1/2}$. This gives
\[
 C_* =3\beta-\beta^2/H-27J_y^2/2
 <3(1.52)-(1.50)^2/4.39-27(0.33)^2/2<2.6 .
\]

Write $r=t/\sqrt3$ and first take $1<r\le7/4$. Then $z=c/r^2$ and
$d=c-z>0$. By the arctangent subtraction formula and $\arctan u\le u$ for
$u\ge0$, the value of the tangent in \eqref{eq:fulltangent} at $a=a_0$,
divided by $t$, is at most
\[
 T=H/r+\frac{27}{2}\frac{r-1}{r(1+3r)} .
\]
We show $G(z)>T$. On $e^\perp$, the identity
$Q_z=r^{-2}Q_c+(1-r^{-2})A$ implies
\[
 (Q_z-T)(v,v)\ge D Q_c(v),\qquad
 D=\frac{9+55/r^2}{64}-\frac{T}{4H}.
\]
Here $T\ge4H/7>\ell$ and $Q_c\ge4H\|\cdot\|_2^2$. On $e$ the form
$Q_z-T$ equals $H-d\beta-T$, and its pairing with $v\perp e$ is
$-dB(e,v)$. By \eqref{eq:stream}, completing the square yields the desired
strict lower bound if $D>0$ and $(H-d\beta-T)D/d^2>C_*$.

Substitution and factorization give
\begin{align*}
 \frac{(H-d\beta-T)D}{d^2}
  &=\frac9{16}\frac{8H+27(3r-1)/(3r+1)}{(r+1)^2}\\
  &\quad\cdot
 \left[\frac{55+9r^2}{64}-\frac r4-\frac{27r(r-1)}{8H(1+3r)}\right].
\end{align*}
The first factor exceeds $3.96$. Its numerator has logarithmic derivative
less than $2/(r+1)$, since $162(r+1)/(3r+1)^2\le2\cdot162/16<2(8H)$.
Thus $r=7/4$ and $8H>35$ give the lower bound
$9(35+27\cdot17/25)/121$. The bracket is $r^2D$. With $y=r-1$ and
$r/(1+3r)\le7/25$, it is bounded below by
\[
  \frac34+\frac y{32}+\frac9{64}y^2-0.22y>\frac23,
\]
as follows by completing the square. In particular $D>0$ and the product
exceeds $3.96(2/3)>2.6$. This proves strict support up to $r=7/4$.

Finally take $t_*=7\sqrt3/4$, $z_*=t_*^{-2}$ and
$T_*=4H/7+162/175$, the value of $T$ at the joining point. For $0<z\le z_*$
concavity and the horizontal Poincar\'e bound give
$G(z)\ge\ell+(T_*-\ell)z/z_*$, using the right limit at zero. Thus
\[
 F(2\arctan t)\ge \ell t+(T_*-\ell)t_*^2/t .
\]
At $t=t_*$ this bound already exceeds the tangent. Its $t$-derivative is
increasing and starts at $2\ell-T_*>1.49$, while the $t$-derivative of the
tangent is $(27/2)/(1+t^2)\le216/163<1.49$. The tangent therefore stays strictly below $F(\alpha)$ for every
$a_0<\alpha<\pi$. For $a<a_0$ strict convexity up to
$a_0$ gives both a smaller tangent value at $a_0$ and a smaller slope.
This proves Proposition~\ref{prop:support}.

\begin{proposition}\label{prop:convex}
If $P$ is convex and $n(P)\le N$, then $J(P)\ge r_N$. Equality holds
precisely for polygons similar to $R_N$.
\end{proposition}
\begin{proof}
Take $a=2\pi/N$. The tangent at $a$ is strictly below $F(0)=0$ as well. Indeed, the difference is positive at $a/2$ and increases towards zero opening.
If there are fewer than $N$ cells, append zeros to the list of openings until it has length $N$. Their sum is $Na$,
so the linear terms in the supporting-line inequalities cancel. Equation
\eqref{eq:sumF} gives $J(P)\ge NF(a)=r_N$. Equality requires all $N$
openings to equal $a$. The ground state has positive mass in every cell, so
equality also holds in every mixed-cell comparison. By
Lemma~\ref{lem:cell} and \eqref{eq:cells}, the cells are isosceles triangles
of equal area and equal apex angle. Their bases have equal length. At each
polygon vertex the two adjacent base angles add to $\pi-a$. The polygon is
therefore regular. The regular polygon attains equality by \eqref{eq:sumF}.
\end{proof}

\section{A bound in terms of the number of convex vertices}\label{sec:turns}

\begin{proposition}\label{prop:turns}
Let $m\ge3$ be an integer. If a simple polygon $P$ has at most $m$ convex
vertices, then $J(P)>r_{m+2}$.
\end{proposition}

We use one integral weight for all $m$. The scalar estimates below hold for
every real $m\ge3$.

\subsection{The regular competitor and the torsion deficit}
Normalize both polygons to area $\pi$. Write $t=1/(m+2)\le1/5$,
$j=j_{0,1}$ for the first positive zero of $J_0$, and $z_0=4\zeta(3)$.
The torsional rigidity is $T(P)=\int_P v$, where $v\in H_0^1(P)$ solves
$-\Delta v=1$. Define $D=1-8T(P)/\pi$. The planar Kohler--Jobin inequality
$\lambda_1(P)\sqrt{T(P)}\ge j^2\sqrt{\pi/8}$ gives
$\lambda_1(P)/j^2\ge(1-D)^{-1/2}$ \cite{brasco}.

The Schwarz--Christoffel map of the disk with derivative
$(1-w^{m+2})^{-2t}$ has image a regular polygon of area $\pi Z(t)$, where
\[
 Z(t)=\frac{\Gamma(1+t)^2\Gamma(1-2t)}{\Gamma(1-t)^2\Gamma(1+2t)},\qquad
 \log Z(t)=\sum_{k\ge3,\,k\, {\rm odd}}\frac{(2^{k+1}-4)\zeta(k)}k t^k .
\]
Indeed, integrating along a radius gives circumradius $t\mathrm B(t,1-2t)$,
so $Z(t)=\sin(2\pi t)[t\mathrm B(t,1-2t)]^2/(2\pi t)$. The gamma reflection
formula and the Taylor expansion of $\log\Gamma(1+z)$ give the two
expressions above. The circle mean of the squared modulus of the derivative
is at least one. Transporting the radial disk ground state to the polygon
therefore gives $\lambda_1(R_{m+2})/j^2\le Z(t)$ after scaling to area $\pi$.

Suppose, for a contradiction, that $\lambda_1(P)\le\lambda_1(R_{m+2})$. Then
\begin{equation}\label{eq:dupper}
  D\le1-Z(t)^{-2}\le d=\frac{h}{1+h/2},\qquad h=t^3(2z_0+33t^2).
\end{equation}
In fact, $2\log Z-2z_0t^3\le4(17/16)(2t)^5/[5(1-4t^2)]<33t^5$, and
$1-e^{-h}\le h/(1+h/2)$ follows from $2\operatorname{arctanh}(h/2)\ge h$.
We shall use
\[
 4.8<z_0<4.81,\qquad \zeta(4)<1.083,\qquad \zeta(5)<17/16,\qquad
 \zeta(6)>1+1/64 .
\]
For $\zeta(3)$ a lower bound is the first six terms plus
$\int_7^\infty x^{-3}\dd x$. An upper bound is the first four terms plus
$\int_{9/2}^\infty x^{-3}\dd x$. For $\zeta(4)$ take three terms and the
midpoint integral starting at $7/2$. For $\zeta(5)$ take two terms and the
tail integral starting at two.

\subsection{Conformal normalization and the torsion identity}\label{sec:torsiontool}
Choose $f:\D\to P$ centered at a point of maximum conformal radius.
Translate the image and rotate the parameter so that $f(0)=0$ and
$a=f'(0)>0$. To justify this choice, start with any Riemann map $f_0$.
Its derivative belongs to the Bergman space by the area formula.
Cauchy--Schwarz on the Taylor coefficients gives
$(1-|z|^2)|p(z)|\le(\pi^{-1}\int_\D|p|^2)^{1/2}$ for every Bergman function
$p$. Polynomial approximation implies that $(1-|z|^2)|f_0'(z)|$ tends to
zero as $|z|\to1$. Move its maximizer to the origin by a disk automorphism.
Stationarity at the origin gives $f''(0)=0$. Write
\[
 f(z)=\sum_{j\ge1}a_jz^j,\quad a_1=a,\ a_2=0,\quad
 b=a^2,\ B=1-b,\quad f'/a=1+\sum_{k\ge2}g_k z^k.
\]
The area identity reads $\sum j|a_j|^2=1$. Introduce
\[
 \begin{gathered}
 S_k=|g_k|^2/(k+1),\quad S=\sum S_k=B/b,\quad Q=\sum\frac{k-1}{k+1}S_k,\\
 V=S-Q,\qquad Q^\sharp=\sum_{k\ge4}\frac{k-3}{(k+1)^2}|g_k|^2 .
 \end{gathered}
\]
We shall use three consequences of Wang's torsion identity.

\begin{lemma}\label{lem:torsion}
With the preceding normalization,
\begin{equation}\label{eq:coercive}
 D\ge B^2/4,\quad D\ge2bB/3,\quad
 D/b^2\ge2Q+\tfrac32V^2-2V\sqrt{Q-V/2},\qquad 0\le V\le2Q .
\end{equation}
\end{lemma}

\begin{proof}
For $s\ge2$ and $1\le p,q<s$, put $A_{p,s}=a_p a_{s-p}$ and
$K_{p,q,s}=\min(p,q,s-p,s-q)$. Then
\begin{equation}\label{eq:sos}
 D=\tfrac12\sum_{s\ge2}\sum_{p,q=1}^{s-1}K_{p,q,s}|A_{p,s}-A_{q,s}|^2 .
\end{equation}
To derive the identity, first assume $f$ extends analytically beyond the
unit circle. If $\mathcal H$ denotes harmonic extension from the circle,
the pullback of the torsion function is $(\mathcal H(|f|^2)-|f|^2)/4$.
Energy orthogonality expresses $16T(P)$ as the energy of $|f|^2$ minus
that of $\mathcal H(|f|^2)$. These energies are
$\pi\sum_s s|\sum_p A_{p,s}|^2$ and
$4\pi\sum_{l\ge1}l|\sum_{j\ge1}a_{j+l}\bar a_j|^2$, respectively.
Collecting coefficients according to $s$ and $p-q$ gives $s/2-|p-q|$ as
the matrix entry for $8T(P)/\pi$. Symmetrizing in $q,s-q$ gives $K_{p,q,s}$.
For $r=\min(p,s-p)$ its row sum is
$2\sum_{q<r}q+(s-2r+1)r=p(s-p)$. Subtraction of the resulting quadratic
form from $(\sum j|a_j|^2)^2$ yields \eqref{eq:sos}. For a general map
apply the calculation to $f(uz)$ for $0<u<1$, using the square of its
normalized area in place of one. As $u\uparrow1$ the nonnegative summands
increase, and the torsional rigidities converge by the variational principle
and exhaustion by compactly supported tests.

Let $M$ be the maximum of the torsion function, and write $A(t)=|\{v>t\}|$.
The flux identity and coarea formula, followed by Cauchy--Schwarz and the
planar isoperimetric inequality, give $-A'(t)\ge4\pi$ for almost every
$t<M$. Integration yields $T(P)\le\pi M-2\pi M^2$, hence
$D\ge(1-4M)^2$. Recenter the map at the torsion maximum and translate its
image to put this point at zero. If the new coefficients are $\tilde a_j$,
the harmonic extension formula at zero gives $4M=\sum|\tilde a_j|^2$.
Since $\sum j|\tilde a_j|^2=1$ and $|\tilde a_1|^2\le b$, it follows that
$4M\le(1+b)/2$. This proves $D\ge B^2/4$.

Put $Y=\sum_{j\ge3}|a_j|^2=bV/2$. In the area square minus the kernel form,
the endpoint terms $p=1,s-1$ contribute $2(s-3)b|a_{s-1}|^2$ for $s\ge4$,
with sum $2b^2Q$. The remaining nonzero products have $3\le p\le s-3$.
Bound the interior quadratic form by its row sums. Each such row loses
six upon deletion of $q=1,2,s-2,s-1$, leaving at least $6Y^2$.
On $s=k+2$ there are at most $k-3$ interior terms. Cauchy--Schwarz bounds
the absolute endpoint--interior contribution by
\[
 4a\Big(\sum_{k\ge4}(k-3)|a_{k+1}|^2\Big)^{1/2}Y
 =4b\sqrt{Q^\sharp}\,Y .
\]
Thus $D/b^2\ge2Q+3V^2/2-2V\sqrt{Q^\sharp}$. Termwise comparison gives
$Q^\sharp\le(3Q-S)/2=Q-V/2$ and $Q\ge S/3$. Moreover,
$2V\sqrt{Q^\sharp}\le3V^2/2+2Q^\sharp/3\le3V^2/2+Q-S/3$.
It follows that $D/b^2\ge2S/3$, which gives the second bound. Absolute
convergence of the kernel forms follows from their row sums and
$2|uv|\le|u|^2+|v|^2$.
\end{proof}

\subsection{A shifted moment inequality and its integral}
Let $\xi_i\in\partial\D$ be the prevertex with interior angle $\omega_i$.
Put $\nu_i=(\pi-\omega_i)/(2\pi)$. These weights sum to one, and the
number $p$ of positive weights is the number of convex vertices. Thus
$p\le m$. The Schwarz--Christoffel formula gives
\[
 h=f'/a=\prod_i(1-\bar\xi_i z)^{-2\nu_i},\qquad
 H(z)=\sum_i\frac{\nu_i}{1-\bar\xi_i z}-1,\qquad 2hH=zh'.
\]
For $c\ge0$ the rational function $B_c=(H-c)/(1+c+H)$ has modulus one on
the circle, since $\Re H=-1/2$ away from the prevertices. Its degree is
the number $n$ of effective vertices. Its preimages of one are exactly the
$n$ prevertices, all simple, with $\xi_iB_c'(\xi_i)=(1+2c)/\nu_i$.
Counting signed crossings of one gives winding number $2p-n$.
Reflection in the unit circle exchanges zeros and poles, so their total
number in the disk is $n$. The argument principle therefore gives exactly
$p$ zeros in the disk, counted with multiplicities. Write $B_c$ as a
unimodular constant times a quotient of finite Blaschke products.
On a circle of radius $r$ every numerator factor has mean log modulus at
least $\log r$, and each denominator factor has modulus at most one.
Jensen's formula for $1+c+H$ gives mean log modulus at least $\log(1+c)$.
Also $h$ has no zeros and $h(0)=1$, so its mean log modulus is zero.
The arithmetic--geometric mean inequality now yields
\[
 \int_{\partial\D}|h(r\zeta)(H(r\zeta)-c)|^2\frac{|\dd\zeta|}{2\pi}
       \ge (1+c)^2 r^{2p}.
\]
By Parseval and $2hH=zh'$, the left side is
$c^2+\sum_{k\ge2}(k/2-c)^2|g_k|^2r^{2k}$. Put $r=e^{-s/2}$,
$M(s)=(e^{ms}-1)^{-1}$, and $c=M(s)$. Since $r^{2p}\ge e^{-ms}$ and
$(1+M)^2e^{-ms}-M^2=M$, we obtain
\begin{equation}\label{eq:jensen}
 4 M(s)\le \sum_{k\ge2}|g_k|^2e^{-ks}\big(k-2M(s)\big)^2 .
\end{equation}
Use the positive weight and its Laplace transform
\[
 w(s)=\tfrac12s^2(1+s)e^{-2s},\qquad
 L(x)=\int_0^\infty w(s)e^{-xs}\dd s=(x+2)^{-3}+3(x+2)^{-4}.
\]
For each $k\ge2$ we claim
\begin{equation}\label{eq:nonpositive}
 \int_0^\infty w(s)e^{-ks}(M^2-kM)\dd s
       =\sum_{j\ge1}(j-1-k)L(k+mj)\le0.
\end{equation}
To see the sign, set $q=1+(k+2)/m$ and put probability weights proportional
to $(q+i)^{-3}$ on $i=0,1,\ldots$. A midpoint upper bound for
$\sum(q+i)^{-2}$ and a trapezoidal lower bound for $\sum(q+i)^{-3}$ give
\[
 \mathbb E i\le
 \frac{1/(q-1/2)}{1/(2q^2)+1/(2q^3)}-q
 =q-\frac{q(q-1)}{(q-1/2)(q+1)} .
\]
If $q\le2$, the right side is at most two. If $q\ge2$, the subtracted
fraction is at least $1/3$, and $q-1/3\le(k+4)/3\le k$. Hence
$\mathbb E i\le k$. Reweighting by the decreasing function $(q+i)^{-1}$
only decreases the mean. Taking $i=j-1$ proves the claim for both powers
in $L$. For fixed $k$ the series are absolutely convergent.

Finally $k^2(k+1)L(k)\le(k-1)/(k+1)$. On clearing positive denominators
the difference has numerator $5(k-2)^3+33(k-2)^2+56(k-2)+4$. Integrate
\eqref{eq:jensen} against $w$. Tonelli's theorem applies to its square
terms, and \eqref{eq:nonpositive} bounds their cross contributions mode by
mode. We conclude that
\begin{equation}\label{eq:rawlower}
                     Q\ge4\sum_{j\ge1} L(mj).
\end{equation}

\subsection{A uniform separation of the two bounds}
Put $x_j=2t(1-1/j)$, so $mj+2=(j/t)(1-x_j)$. On the right of
\eqref{eq:rawlower} expand the inverse cube through $x_j^3$ and the inverse
fourth power through $x_j^2$. The omitted terms are nonnegative. This gives
\[
 \frac Q{t^3}\ge z_0+(6z_0-12\zeta(4))t+
 96(\zeta(3)-\zeta(4))t^2+
 4(80\zeta(3)-120\zeta(4)+40\zeta(6))t^3 .
\]
Consequently $Q\ge a=t^3 A(t)$ with
$A(t)=z_0+15.8t+11.2t^2+26t^3<8.63$.

Equation \eqref{eq:dupper} implies $d<0.084$ and $d/t^2<2.1$. Indeed
$h/t^3<10.94$, and $10.94t/(1+10.94t^3/2)$ increases up to $t=1/5$,
where it is less than $2.1$. By \eqref{eq:coercive}, $B<0.58$.
For $0.15\le B\le.58$ one has $2B(1-B)/3\ge.085$, so $b>0.85$. Hence
\[
 b^{-2}-1=(1+b)B/b^2\le\frac{3(1+b)}{2b^3}D<4.52d .
\]
Since $h-d=hd/2\ge d^2/2$, this gives
\begin{equation}\label{eq:upfinal}
                 D/b^2\le d+4.52d^2\le h+4.02(2.1)^2t^4<h+17.73t^4 .
\end{equation}

For a lower bound, first observe that \eqref{eq:rawlower} gives $Q>0$.
Concavity of the square
root yields $\sqrt{Q-V/2}\le\sqrt Q-V/(4\sqrt Q)$. Completing the square
in $V$ in \eqref{eq:coercive} gives
\[
 \frac{D}{b^2}\ge 2Q-\frac{2Q^{3/2}}{1+3\sqrt Q}
        \ge 2a-\frac{2a^{3/2}}{1+3\sqrt a}.
\]
The last inequality uses that $2Q(1+2\sqrt Q)/(1+3\sqrt Q)$ increases in $Q$.
Now
\[
 (2a-h)/t^4=31.6-10.6t+52t^2\ge31.6-10.6^2/208>31
\]
and $\sup_{y>0} y/(1+3\sqrt A\,y^3)=2/[3(6\sqrt A)^{1/3}]$. Thus
\[
 \frac{2a^{3/2}}{(1+3\sqrt a)t^4}
   =\frac{2A^{3/2}\sqrt t}{1+3\sqrt A(\sqrt t)^3}
   \le\frac{4 A^{4/3}}{3\sqrt[3]6}<13.1 .
\]
It follows that $D/b^2>h+17.9t^4$, contradicting \eqref{eq:upfinal} and
proving Proposition~\ref{prop:turns}.

\section{Compactness with one reflex vertex}\label{sec:compact}

\begin{lemma}\label{lem:compact}
Let $P_i$ be simple polygons of area $A>0$, each with at most $N$ effective
sides and exactly one reflex vertex. If $\liminf_i\lambda_1(P_i)<\infty$,
there is a simple polygon $Q$ such that
\begin{equation}\label{eq:compact}
  0<|Q|\le A,\qquad n(Q)\le N,\qquad
  \lambda_1(Q)\le\liminf_i\lambda_1(P_i).
\end{equation}
\end{lemma}

\begin{proof}
Take a subsequence realizing the lower limit and having a fixed side count
$p\le N$. Translate its reflex vertex $v_i$ to the origin. A ray from $v_i$
splits the interior angle $\omega_i>\pi$ into two angles smaller than $\pi$
if its direction, measured from one side, lies in $(\omega_i-\pi,\pi)$.
Choose such a ray avoiding the other vertices and all directions parallel
to edges. Its first return to the boundary is a point $z_i$ in the relative
interior of a nonincident edge $E_i$. The chord $[v_i,z_i]$ cuts the polygon
into two pieces with all interior angles less than $\pi$. Their closures
$K_{1i},K_{2i}$ are convex and lie on opposite sides of the chord line.
Each of the two boundary paths from $v_i$ to $z_i$ contains another original
vertex, so each piece has at most $p$ vertices. Both pieces retain a
subsegment of $E_i$ with the same induced orientation. The line of $E_i$
therefore supports both pieces on the same side (Figure~\ref{fig:cut}).

\begin{figure}[htbp]
\centering
\begin{tikzpicture}[x=0.73cm,y=0.67cm,baseline=(current bounding box.center),
                    every node/.style={font=\small},line join=round]
 \fill[ink!15] (-2,0)--(0,0)--(0,1.12)--(-2,1.93)--cycle;
 \fill[ink!5] (0,0)--(2,0)--(2,1.93)--(0,1.12)--cycle;
 \draw[thick] (-2,0)--(2,0)--(2,1.93)--(0,1.12)--(-2,1.93)--cycle;
 \draw[ink,densely dashed,thick] (0,0)--(0,1.12);
 \fill (0,0) circle (1.45pt) (0,1.12) circle (1.45pt);
 \node[above] at (0,1.12) {$v_i$};
 \node[below] at (0,0) {$z_i$}; \node[below] at (-1,0) {$E_i$};
 \node at (-1.04,0.73) {$K_{1i}$}; \node at (1.04,0.73) {$K_{2i}$};
 \begin{scope}[shift={(6.8,0.91)}]
  \fill[ink!15] (-2,0)--(-1.65,0.92)--(0.15,1.02)--(0.74,0)--cycle;
  \fill[ink!5] (-0.62,0)--(0.16,-1.06)--(1.64,-0.9)--(2,0)--cycle;
  \draw[thick] (-2,0)--(-1.65,0.92)--(0.15,1.02)--(0.74,0)--(2,0)--(1.64,-0.9)--(0.16,-1.06)--(-0.62,0)--cycle;
  \draw[ink,densely dashed,thick] (-0.62,0)--(0.74,0);
  \node at (-0.55,0.48) {$K_1$}; \node at (0.66,-0.5) {$K_2$};
  \node[right] at (2.05,0) {$L$};
 \end{scope}
\end{tikzpicture}
\caption{On the left the chord from the reflex vertex ends in an original
edge. On the right two limiting pieces meet along a segment of positive
length. The relative interior of the dashed overlap lies in the union's
interior.}\label{fig:cut}
\end{figure}

Extend the $L^2$-normalized ground states by zero to $u_i\in H^1(\R^2)$.
Their Dirichlet energies are bounded. A convex piece of area at most $A$
and diameter $h$ lies in a strip of width at most $2A/h$. Indeed, form two
triangles on a diameter by taking a farthest point on each side of its line.
The sum of their areas is half the diameter times the width in the
perpendicular direction. For any strip $L_w$ of width $w$,
\[
 \int_{L_w}|u|^2\le 2w\|u\|_{2}\|\nabla u\|_{2}\qquad(u\in H^1(\R^2)).
\]
Apply the one-dimensional bound $\sup|g|^2\le2\|g\|_2\|g'\|_2$ normally
to the strip, then integrate tangentially. Thus any piece of diverging
diameter carries vanishing $L^2$ mass. Pass to a further subsequence along
which both diameters have limits in $[0,\infty]$. At least one is bounded.

If the other diameter diverges, the bounded piece lies in a fixed ball since
it contains the origin. Its vertices have a subsequence converging to a
convex hull $K_1$ with at most $N$ vertices. Take a weak $H^1(\R^2)$ limit
$u$ and strong $L^2$ limits on bounded balls. Vanishing mass on the
unbounded piece gives $\|u\|_2=1$ and $\operatorname{supp}u\subset K_1$.
In particular $K_1$ has positive area. The zero-extension trace criterion
gives $u\in H_0^1(K_1^\circ)$. Weak lower semicontinuity of the energy and
continuity of convex area give \eqref{eq:compact} with $Q=K_1^\circ$.

Suppose now that both diameters are bounded. Keep the index of $E_i$ fixed
and take limits of all original vertices, $z_i$, and oriented unit normals
to the chord and to $E_i$. Let $\mathcal V$ be the set of limits of the $p$
original vertices and put $z=\lim_i z_i$. The Hausdorff limits $K_1,K_2$
of the two pieces are convex hulls of at most $p$ points each, all in
$\mathcal V\cup\{z\}$. They lie on opposite sides of the limiting chord
line $L$ and on the same side of the limiting support line of $E_i$.
Continuity of area and disjointness of the interiors give
$|K_1|+|K_2|=A$.

Put $U=\operatorname{int}(K_1\cup K_2)$. If one piece degenerates to a
point or a segment, $U$ is the interior of the other. If both have positive
area and meet only at a point, $U$ is the disjoint union of their interiors.
In either case each component has at most $N$ sides. In the remaining case,
the intervals $I_a=K_a\cap L$, $a=1,2$, overlap in positive length.
Removing the relative interior of $I_a$ from each boundary leaves two
outer polygonal arcs, one on each side of $L$. Join the left endpoints and
the right endpoints along $L$, omitting any zero joins. These two joins are
disjoint, and the outer arcs leave $L$ except at their endpoints. The
result is a polygonal Jordan curve enclosing exactly $U$.

Every turn of its boundary is a vertex of a limiting piece. If
$z\notin\mathcal V$, the endpoints of $E_i$ converge to opposite sides of
$z$ along its limiting line. The two pieces retain the two subsegments of
this limiting edge. The support line and $L$ are distinct, since coincidence
would force one of the two pieces to degenerate. Join a point of the common
overlap near $z$ to the limiting endpoints of the edge. Each half of the
resulting triangle belongs to the corresponding convex piece. Hence $z$
lies in the relative interior of a straight side of $U$. All effective
vertices belong to $\mathcal V$, so $n(U)\le N$.

Finally, Rellich compactness on a fixed ball gives a limit $u$ of unit
$L^2$ norm and energy at most $\liminf_i\lambda_1(P_i)$. It is supported
on $K_1\cup K_2$, which differs from $U$ by a set of area zero.
When $U$ is a single polygon its Lipschitz trace gives $u|_U\in H_0^1(U)$.
When the two open components touch at just one boundary point, their
traces vanish off the contact point and hence almost everywhere on each
boundary. In this case the restriction to each component also has zero
trace. Choose a component $Q$ of lowest first eigenvalue. Then
\[
 \lambda_1(Q)=\lambda_1(U)\le\int_{\R^2}|\nabla u|^2\le\liminf_i\lambda_1(P_i),
 \qquad 0<|Q|\le A.
\]
The side count established above proves the lemma.
\end{proof}

\section{Stationarity and the boundary pressure}\label{sec:pressure}
Let $u>0$ be a first eigenfunction on a simple polygon $P$, and write
$\lambda=\lambda_1(P)$, $n=n(P)$. Define the root-mean-square height and the
pressure on the open sides by
\[
     t_*^2=|P|^{-1}\int_Pu^2,\qquad q=(\partial_\nu u)^2
\]
Following \cite{cheng}, we show that if a nonconvex polygon is stationary
under variations of its supporting lines, then $\{u>t\}$ has at least two
components for almost every $t$ in some neighborhood of $t_*$.

\subsection{Corner expansions and two moments on each edge}
Near a vertex of angle $\omega$ use polar coordinates with $0<\theta<\omega$.
Set $\beta=\pi/\omega>1/2$ and $\gamma=\min(\beta,2)$. Then
\[
 u=c r^\beta\sin(\beta\theta)+R,\quad c>0,\qquad
 |\nabla^j R|=O(r^{\beta+\gamma-j})\quad(j=0,1).
\]
Indeed, take $r_1$ small enough that $\sqrt\lambda r_1<j_{\beta,1}$.
Expanding in angular sines and solving the finite-energy radial Bessel
equations gives
$u=\sum_{m\ge1}c_mJ_{m\beta}(\sqrt\lambda r)\sin(m\beta\theta)$.
The first sine coefficient on $r=r_1$ is positive. Parseval's identity
and $J_\nu(s)=(s/2)^\nu[1+O(s^2/(\nu+1))]/\Gamma(\nu+1)$ give uniform
convergence of the higher modes and their first derivatives on $r\le r_1/2$.
They are $O(r^{2\beta})$ and $O(r^{2\beta-1})$, respectively. On both sides
of the corner,
$q=c^2\beta^2r^{2\beta-2}(1+O(r^\gamma))$. This is integrable. It tends to
zero at convex vertices and to infinity at reflex vertices.

Translate and rotate the supporting line of one effective side, keeping the
others fixed. These variations realize every affine normal velocity
$h(s)=d_0+d_1s$ on that side, where $0<s<L$ is arclength. Consecutive lines
meet transversely, so small variations preserve simplicity and side count.
Pull back the Rayleigh quotient by piecewise affine interpolation on a
fixed triangulation. Simplicity of the first eigenvalue gives differentiability
by analytic perturbation \cite{kato}. The stress tensor in the volume
derivative is $\mathcal T=(|\nabla u|^2-\lambda u^2)I-2\nabla u\otimes\nabla u$.
It is divergence free and has traction $\mathcal T\nu=-q\nu$ on each open
side. Excision of radius-$\epsilon$ sectors at the vertices gives circular
flux $O(\epsilon^{2\beta-1})\to0$. Adding the area derivative yields
\[
 J'(0)=-\frac{|P|}{\int_P u^2}\int_0^L(q-\lambda t_*^2)h(s)\dd s .
\]
Rescaling to restore area leaves $J$ unchanged. Thus at a global minimizer
both affine variations on each side are stationary and
\begin{equation}\label{eq:momentsE}
       \int_0^L(q-\lambda t_*^2)\dd s=\int_0^L s(q-\lambda t_*^2)\dd s=0 .
\end{equation}
Odd reflection gives analyticity of $q-\lambda t_*^2$ on each open side.
The endpoint limits imply it has finitely many zeros. If it had at most one
sign change, a constant or a linear function vanishing at that change would
contradict \eqref{eq:momentsE}. A side joining two vertices of the same type
therefore has at least two sign changes. A side joining different types has
at least three, by parity. For a nonconvex polygon the number $k$ of sides
of the latter kind is even and at least two. Choose points with opposite
strict signs on the two sides of each change. These crossings persist for
all nearby heights. Writing $\mathcal V_P$ for the polygon vertices, we obtain at regular heights near $t_*$
\begin{equation}\label{eq:countlow}
 Z(t)=\#\{x\in\partial P\setminus\mathcal V_P:q(x)=\lambda t^2\}\ge 2n+k .
\end{equation}

\subsection{Flux of the gradient direction and the number of components}
Let $Q=|\nabla u|^2$, $\Psi=(Q+\lambda u^2)/2$, and $\Sigma=\{\nabla u=0\}$.
Off $\Sigma$ define $X=\nabla\Psi/Q$. If $\vartheta$ is a local argument
of $\nabla u$ and $R$ is counterclockwise rotation through $\pi/2$, the
eigenfunction equation gives
\begin{equation}\label{eq:flux-identity}
 \nabla\Psi=-\operatorname{cof}(D^2u)\nabla u,\qquad
 X=R\nabla\vartheta,\qquad \operatorname{div}X=0 .
\end{equation}
Odd reflection of $u$ makes $\Psi$ even across open sides. The Hopf lemma
gives $Q>0$ there, so the normal flux of $X$ vanishes. The corner expansion
also gives $|\nabla u|=c\beta r^{\beta-1}(1+O(r^\gamma))$, uniformly in
the angular variable. Thus $\Sigma$ is contained in a compact subset of $P$.

In fact $\Sigma$ is finite. At any accumulation point,
$\operatorname{tr}D^2u=-\lambda u<0$, so the zero set of at least one first
partial derivative is a regular analytic arc. The other partial derivative
has accumulating zeros on the arc and vanishes along it. The same local
description holds at each limit point of such an arc. Any one-dimensional
critical component is therefore a compact regular curve without endpoints
or branches. Its Jordan interior lies in $P$, where the integral of
$\Delta u=-\lambda u$ is strictly negative. On the curve the boundary flux
of $\nabla u$ is zero, a contradiction.

By \eqref{eq:flux-identity}, the outward flux of $X$ on a small circle
around a critical point is $-2\pi$ times its gradient index. At a vertex,
in axes aligned with the sides as in the corner expansion, the leading
gradient direction is
$(\sin((\beta-1)\theta),\cos((\beta-1)\theta))$. Its total change of
argument along the circular arc tends to $\omega-\pi$. This is the limiting
flux for the normal pointing into the excised sector. Excising all critical
points and corners and applying the divergence theorem gives
$\sum_{P\cap\Sigma}\operatorname{ind}(\nabla u)=1$.

Let $C(t)$ be the number of components of $\{u>t\}$. Choose $0<t<\max u$
so that $\lambda t^2/2$ is a regular value of $\Psi$ and $\lambda t^2$ a
regular value of $q$ on each edge. Almost every $t$ qualifies by Sard's
theorem. It is also a regular value of $u$, since $\nabla u=0$ implies
$\nabla\Psi=0$. Each superlevel component is a smooth Jordan domain.
Indeed, the interior of an inner level curve lies in $P$, and strict
superharmonicity forces $u>t$ in that interior. This contradicts the
change of sign across a regular inner level curve. On the boundary of each
component the gradient points inward, so its total index is one.

Write $C_+,C_-$ for the numbers of components of
$\{2\Psi>\lambda t^2\}$ and $\{2\Psi<\lambda t^2\}$. In a high-pressure
component $H$ the sum of the critical point indices equals the number $C_H$
of superlevel components it contains. Indeed, the closure of each superlevel
component lies in the high-pressure set, and $2\Psi=\lambda u^2$ at critical
points. The pressure level avoids all vertices and meets the sides
orthogonally. It is a compact regular one-manifold, consisting of
$b=Z(t)/2$ arcs with boundary endpoints and possibly some closed curves.
The $b$ disjoint arcs cut a topological disk into $b+1$ regions. Thus
$b+1\le C_++C_-$.

The flux identity on a high-pressure component reads
\[
 0<\int_{\partial H\cap P}\frac{|\nabla\Psi|}{Q}\dd s
  =\sum_{v\in\overline H\cap\mathcal V_P}(\omega_v-\pi)+2\pi C_H ,
\]
where the sum is over incident polygon vertices, all reflex. Each such
component must contain a superlevel component or meet a reflex vertex.
If $q_0$ is the total number of reflex vertices, then
$C_+\le q_0+C(t)$. For a low-pressure component $L$, let $a_L$ be its
number of incident convex vertices and $i_L$ its critical point index sum.
Its positive level-boundary flux is
$\sum_{v\in\bar L\cap\mathcal V_P}(\pi-\omega_v)-2\pi i_L>0$. If $a_L=0$ this gives
$i_L\le-1$. If $a_L>0$ it gives $2i_L<a_L$. In both cases
$a_L-i_L\ge1$. Summing and using $\sum_L i_L=1-C(t)$ and
$\sum_La_L=n-q_0$, we obtain $C_-\le n-q_0+C(t)-1$. Combining the bounds
for the two signs gives
\begin{equation}\label{eq:countup}
                        Z(t)\le 2n+4C(t)-4 .
\end{equation}
With \eqref{eq:countlow} this implies $C(t)\ge1+k/4>1$ for almost every
$t$ near $t_*$ on a nonconvex stationary polygon.

\subsection{The pentagon as a common competitor}
Suppose three successive vertices $A,B,C$ form an ear, with the open
diagonal $AC$ in $P$. If $|AC|\le\max(|AB|,|BC|)$, stationarity is
impossible. Relabel so that $|BC|\ge|AC|$. Take coordinates
$A=(0,0)$, $B=(L,0)$, $C=(p,h)$ with $h>0$, so $p\le L/2$. At height $y$
the endpoint abscissas of the ear sum to $L+(2p-L)y/h\le L$. The right
cap $x>L/2$ of the ear therefore reflects into $P$ about $x=L/2$.
The reflection of the open sloping boundary of the cap is interior, using
$(A,C)\subset P$ also when $p=L/2$. The function
$u(L-x,y)-u(x,y)$ has nonnegative boundary values on the cap. Since the
cap's first eigenvalue exceeds $\lambda$, testing with the negative part
and then applying the strong maximum principle shows that the difference
is positive in the cap. The Hopf lemma on $AB$ yields
$q(L-x)>q(x)$ for $L/2<x<L$. This contradicts
$\int_0^L(x-L/2)q(x)\dd x=0$, a consequence of \eqref{eq:momentsE}.

A nonconvex quadrilateral replaces a side $AB$ of its triangular convex
hull by $ARB$, with $R$ in the interior of the hull. Let $C$ be the
third hull vertex and order $A,B$ so that $|CA|\ge|CB|$. Then $CAR$ is
an ear and $|CR|<|CA|$. Consequently a nonconvex stationary polygon must
have at least five effective sides. If it is a global minimizer with side
budget $N$, then $N\ge5$, and the admissible regular pentagon gives
$J(P)\le r_5$.

\section{The pentagonal threshold and projection onto the disk ground state}
\label{sec:disk}
Throughout this section and the next we normalize $|P|=\pi$ and use the disk
estimates of \cite{cheng}. Let $\lambda=\lambda_1(P)$, $j=j_{0,1}$ and
$\mu=j^2$. Write $\davg{v}=\pi^{-1}\int_\D v\dd A$ and
$\norm{v}^2=\davg{|v|^2}$. The positive disk ground state is
$\phi(z)=\Phi(|z|^2)=J_0(j|z|)$. The Bessel norm identity gives
$\Cphi=\norm{\phi}^2=J_1(j)^2$. Write $G$ for the inverse of the Dirichlet
Laplacian on $\D$. Our conventions for its kernel and the dual energy norm are
\[
 G_z(w)=\frac{1}{2\pi}\log\left|\frac{1-\bar zw}{z-w}\right|,\qquad
 GF(z)=\int_\D G_z(w)F(w)\dd A(w),\qquad \|F\|_{-1}^2=\davg{\bar F GF}.
\]
The ground state on $P$ is bounded. Indeed, with $R=\operatorname{diam}P$,
majorize $G_P(x,y)$ by the Green kernel of $B(x,R)$. Cauchy--Schwarz gives
$u(x)\le\lambda R\|u\|_{L^2(P)}/\sqrt{8\pi}$.

\subsection{The disk profile and the regular pentagon}
We use the enclosures
\begin{equation}\label{eq:disk-constants}
  5.783<\mu<5.784,\qquad
  0.2694<\Cphi<0.2697,\qquad
  j_{1,1}^2>14.6.
\end{equation}
Here $j_{\nu,1}$ denotes the first positive zero of $J_\nu$. For a rational
verification, put $\Jcal_n(t)=n!(2/\sqrt t)^n J_n(\sqrt t)$, continued to
$t=0$. Its degree-$d$ Taylor polynomial is
\[
  B_{n,d}(t)=\sum_{i=0}^d\frac{(-t/4)^i n!}{i!(n+i)!}.
\]
One has $4(n+1)\Jcal_n'=-\Jcal_{n+1}$. At each evaluation below the
successive absolute ratios in the omitted tail, $t/[4(i+1)(n+i+1)]$, are
less than one. Odd and even truncations give lower and upper bounds.
With $z=t/4$, the cubic $B_{2,3}=1-z/3+z^2/24-z^3/360$ has $z$-derivative
$-[(z-5)^2+15]/120$, and $B_{2,3}(14.6)>0$. Thus $\Jcal_1$ decreases
strictly on $[0,14.6]$. Since $B_{1,5}(14.6)>0$ we obtain $j_{1,1}^2>14.6$
and strict decrease of $\Jcal_0$ there. Now
$B_{0,5}(5.783)>0>B_{0,6}(5.784)$ bounds $\mu$. Finally
$B_{1,5}(5.784)>0.43168$ and $B_{1,6}(5.783)<0.4318$ give the norm
bounds by $\Cphi=(\mu/4)\Jcal_1(\mu)^2$.

Set $s=r^2$ and $\rho=1-s$. The radial equation $(s\Phi')'=-\mu\Phi/4$
integrates to
\begin{equation}\label{eq:disk-average}
  -\Phi'(s)=\frac{\mu}{4}\int_0^1\Phi(st)\dd t.
\end{equation}
Repeated differentiation gives $(-1)^i\Phi^{(i)}\ge0$. We claim that
\begin{equation}\label{eq:phi-chord-bounds}
  \rho\left(1-\frac{9s}{20}\right)\le\Phi(s)\le p_3(s):=\rho\left(1-\frac{7s}{16}+\frac{s^2}{16}\right).
\end{equation}
The Hermite polynomial at nodes $0,0,1$ is $\rho[1-(\mu/4-1)s]$ and its
remainder is nonnegative. Since $\mu/4-1<9/20$, this proves the lower
bound. At nodes $0,1,1,1$ the cubic is
$d\rho+d\rho^2/2+(1-3d/2)\rho^3$ and its remainder is nonpositive.
Here $d=-\Phi'(1)=jJ_1(j)/2=\Phi''(1)$. In this cubic the coefficient of
$d$ is nonnegative, and $4d^2=\mu\Cphi<25/16$. Substituting $d=5/8$ gives
the upper polynomial $p_3$.

For the regular pentagon take the conformal derivative
$(1-z^5)^{-2/5}=\sum c_mz^{5m}$, where $c_0=1$ and
$c_{m+1}/c_m=(5m+2)/(5m+5)$. The image has area $\pi S$ with
$S=\sum c_m^2/(5m+1)$. Transport $\phi$ to this image as a trial function.
After angular integration, retaining just $m=0,1$ gives a lower bound on
the mass. The second term involves $I=\int_0^1s^5\Phi(s)^2\dd s$, and
\eqref{eq:phi-chord-bounds} gives $I\ge989/336000$. For
$T_m=1/6+1/(10m-2)$ the recursion gives
\[
 T_m-\left(\frac{c_{m+1}}{c_m}\right)^2 T_{m+1}-\frac{1}{5m+1}
 =\frac{6(35m+23)}{25(m+1)^2(5m-1)(5m+1)(5m+4)}>0\quad(m\ge1).
\]
Multiplying by $c_m^2$ and summing from $m=3$ yields
$S\le\sum_{m=0}^2c_m^2/(5m+1)+c_3^2T_3=538286/515625$. Hence
\begin{equation}\label{eq:pent}
 \frac{r_5}{\pi\mu}\le\frac{S}{1+4I/(25\Cphi)}<\frac{5211}{5000}.
\end{equation}
Assume from now on that $J(P)\le(5211/5000)\pi\mu$. The Faber--Krahn
and Kohler--Jobin inequalities yield
\[
 0\le\delta:=1-\mu/\lambda\le\delta_5:=211/5211,\qquad
 D=1-8T(P)/\pi\le2\delta-\delta^2<0.084 .
\]
In the centered conformal coordinates of Section~\ref{sec:torsiontool},
write the derivative $g=f'$ as
\begin{equation}\label{eq:conformal-norm}
  g(z)=a+h(z),\quad a>0,\quad
  h(z)=\sum_{k=2}^\infty b_k z^k,\quad
  e^2:=\norm{h}^2=\sum_{k=2}^\infty|\alpha_k|^2,\quad
  a^2+e^2=1,
\end{equation}
where $\alpha_k=b_k/\sqrt{k+1}$. Put $x=e^2=1-a^2$. By \eqref{eq:coercive},
$x<0.58$. In the interval $[0.15,0.58]$ one would have $2x(1-x)/3\ge.085$.
Thus $x<3/20$. If $x=0$ the domain is a disk and the remainder $W$ below
is identically zero.

\subsection{Holomorphic weights and a mixed Green moment}
For every square-integrable holomorphic $h$ and $Z\in H_0^1(\D)$ we have
\begin{equation}\label{eq:weight}
                 \norm{hZ}\le\norm h\,\norm{\nabla Z}/\sqrt\mu .
\end{equation}
To prove this, take a primitive $H$ of $h$. For any smoothly bounded
$O\Subset\D$, Green's formula expresses $\int_O|h|^2\dd A$ as the sum of
the signed areas enclosed by the images of its boundary curves. Parametrize
each image by constant speed and subtract its mean. Wirtinger's inequality
bounds the absolute signed area by the square of its length divided by
$4\pi$, including for curves with self-intersections. Thus
$4\pi\int_O|h|^2\dd A\le(\int_{\partial O}|h|\dd s)^2$.

Start with a smooth compactly supported test $Z$, approximated by functions
Morse at nonzero heights. For $\nu(t)=\int_{|Z|>t}|h|^2\dd A$, coarea and
Cauchy--Schwarz give
\[
 \int_{|Z|=t}|\nabla Z|\dd s\ge 4\pi\nu(t)/[-\nu'(t)]
\]
for almost every $t$ in the range. On the disk of radius $\norm h$, take
the radial decreasing rearrangement with distribution $\nu$ for ordinary
area. Its energy is at most the original energy, and its squared mass is
$\pi\norm{hZ}^2$. The first eigenvalue of that disk is $\mu/\norm h^2$.
Density gives \eqref{eq:weight}, with the zero weight immediate. Also,
$\big||h_1|^2-|h_2|^2\big|\le|h_1+h_2||h_1-h_2|$. Pairing with $Z$ and
using \eqref{eq:weight} shows that $h\mapsto|h|^2$ is locally Lipschitz
from the Bergman space into $H^{-1}$.

Write $e_4^2=\sum_{k\ge4}|\alpha_k|^2$ and
$I_w=\davg{w|h|^2G(wh)}$ for a smooth radial weight $w=w(s)$. The Poisson
equation in angular mode $l$ gives $G(wz^l)=z^lP_{l,w}(s)$, where
\begin{equation}\label{eq:free-profile}
  P_{l,w}(s)=\frac14\int_s^1\int_0^1 u^l w(tu)\dd u\dd t.
\end{equation}
For polynomial $h$ set $n=k+l$ and $d_{kl}=\sqrt{(k+1)(l+1)(n+1)}$.
Angular orthogonality gives
\[
  I_w=\sum_{k,l\ge2}m_{kl}\alpha_k\alpha_l\overline{\alpha_{k+l}},
  \qquad
  m_{kl}=d_{kl}\int_0^1 s^n w(s)P_{l,w}(s)\dd s.
\]
Cauchy--Schwarz along $k+l=n$ and then in $n$ gives
$|I_w|\le e^2e_4\sup_{n\ge4}(\sum_{k+l=n}|m_{kl}|^2)^{1/2}$.
For $w=\Phi$ we have $\Phi\le\rho(1-3s/8)$ and
$P_{l,\Phi}\le P_{2,\rho}=\rho(2+3\rho)/96$. In a beta integral
$\int_0^1s^n\rho^j\dd s$ for $j=2,3$, inserting $1-3s/8$ changes the
integral by a factor $1-3(n+1)/[8(n+j+2)]\le19/24$. Since
$(k+1)(l+1)\le(n+2)^2/4$, this gives
\[
  \left(\sum_{\substack{k,l\ge2\\ k+l=n}}|m_{kl}|^2\right)^{1/2}
    \le\frac{19}{24}\,\frac{2+9/(n+4)}{96(n+3)}\sqrt{\frac{n-3}{n+1}}
    \le\frac{19}{24}\,\frac{25}{768\sqrt{180}}<\frac{1}{500},
\]
For the last comparison put $v=n-3$ and use
$(n+3)^2(n+1)/(n-3)=v^2+16v+84+144/v\ge180$.
Pass to polynomial limits using \eqref{eq:weight} and boundedness of
multiplication by a smooth weight on $H_0^1$ and its dual. We obtain
\begin{equation}\label{eq:mixedI}
                         |I_\Phi|\le e^2e_4/500 .
\end{equation}

\subsection{Energy and mass of the orthogonal remainder}
Put $p=|a+h|^2$, $\kappa=\lambda a^2$, and
$c=\davg{(u\circ f)\phi}/\Cphi>0$. Normalize the pullback by writing
\begin{equation}\label{eq:V-pde}
  V=\frac{u\circ f}{c}=\phi+W,\qquad W\perp\phi,\qquad -\Delta V=\lambda pV.
\end{equation}
Testing with $\phi\circ f^{-1}$ yields
$\lambda\le\mu\Cphi/\davg{p\phi^2}\le\mu/a^2$. In particular
$\delta\le x$ and $\kappa\le\mu$. Set $d=1/2-x$, $w=\norm{\phi h}^2$
and $F=\phi(2a\Re h+|h|^2)$. We claim
\begin{align}
  \left(d+\frac{w}{\Cphi}\right)\bigl[(x-\delta)\Cphi-w\bigr]
    &\le\mu\|F\|_{-1}^2-\frac{w^2}{\Cphi},
    \label{eq:projected-residual}\\
  \left(d+\frac{w}{\Cphi}\right)\norm{\nabla W}
    &\le\mu\|F\|_{-1}.
    \label{eq:W-energy}
\end{align}
Define $B(Y,Z)=\lambda^{-1}\davg{\nabla Y\cdot\nabla Z}-\davg{pYZ}$.
In the energy inner product, the compact positive operator $\mu G(p\,\cdot)$
has eigenvalues $\mu/\lambda_i(P)$. At least one nodal domain of a second
eigenfunction has area at most $\pi/2$. Faber--Krahn on that domain yields
$\lambda_2(P)\ge2\mu$. If $Z\perp\phi$, the trace of the compression to
$\operatorname{span}\{\phi,Z\}$ is at most $\mu/\lambda+1/2$. The diagonal
entry in the energy-unit direction of $\phi$ equals $1-x+w/\Cphi$.
Subtracting it gives $\mu B(Z,Z)\ge(d+w/\Cphi)\norm{\nabla Z}^2$.
By the equation, $B(V,\cdot)=0$ and
$B(W,W)=B(\phi,\phi)=(x-\delta)\Cphi-w$. Moreover,
$B(W,Z)=\davg{FZ}$ for $Z\perp\phi$. Since $\davg{F\phi}=w$ and
$G\phi=\phi/\mu$, the projected source $F_\perp=F-w\phi/\Cphi$ satisfies
$\mu\|F_\perp\|_{-1}^2=\mu\|F\|_{-1}^2-w^2/\Cphi$. Estimating
$B(W,W)=\davg{F_\perp W}$ by duality proves both inequalities.

\begin{proposition}\label{prop:disk}
Under the preceding spectral bound, set $\eta=1-\kappa/\mu$. Then
\begin{equation}\label{eq:energy}
  e<\frac27,\qquad
  \norm{\nabla W}<0.29,\qquad
  \norm{W}<0.076,\qquad
  0\le\eta\le\frac35 e^2<\frac{1}{20}.
\end{equation}
\end{proposition}
\begin{proof}
It suffices to treat $x>0$. Put $y=e_4/e$, $H=w/x$ and
$T=\|\phi\Re h\|_{-1}^2/x$. Use the constants
\[
  M=5.784,\quad m=0.0644,\quad q=0.0372,\quad t=0.001202,\quad v=0.000976,
  \quad c_1=\frac{M}{125}.
\]
The normalized modal coefficients of $H$ and $T$ for $k\ge2$ are
\[
  (k+1)\int_0^1 s^k\Phi(s)^2\dd s,\qquad
  \frac{k+1}{2}\int_0^1 s^k\Phi(s)P_{k,\Phi}(s)\dd s.
\]
As $k$ increases, the probability measures $(k+1)s^k\dd s$ shift toward
one. Both profiles decrease with $s$, and $P_{k,\Phi}$ also decreases with
$k$. It suffices to bound the integrals at $k=2,4$ with $p_3$ in place of
$\Phi$. For $H$ the two polynomial integrals give $173/2688,2407/88704$.
For $T$ they give $30991/25804800,8600083/38149816320$. Hence
$H\le m-qy^2$ and $T\le t-vy^2$. By \eqref{eq:weight},
$\|\phi|h|^2\|_{-1}\le e\norm{\phi h}/\sqrt\mu$. The cross term in
$\|F\|_{-1}^2$ is $4a\Re I_\Phi$, so \eqref{eq:mixedI} gives
\begin{equation}\label{eq:raw-residual}
  \frac{\mu}{x}\|F\|_{-1}^2\le 4M(1-x)(t-vy^2)+xH+c_1\sqrt x\,y.
\end{equation}
On expanding \eqref{eq:projected-residual}, the terms $-w^2/\Cphi$ cancel.
After division by $x$ the coefficient of $H$ on the right is $d+\delta$.
Since $x<3/20$, the coefficient of $-y^2$ is at least
$h_0=7q/20+(17/20)4Mv$. Completing the square in $y$ gives
\begin{equation}\label{eq:scalar}
  d\left(1-\frac{\delta}{x}\right)\Cphi
    \le m(d+\delta)+4Mt(1-x)+\frac{c_1^2 x}{4h_0}.
\end{equation}
Set $X=4/49$, $d_X=1/2-X$, $C_0=0.2694$ and $C_1=0.2697$. The concave
quadratic
\[
  (x-\delta_5)\left(\frac12-x\right)C_0
    -x\left[m\left(\frac12-x+\delta_5\right)+4Mt(1-x)+\frac{c_1^2 x}{4h_0}\right]
\]
exceeds $1/40000$ at both endpoints of $[X,3/20]$. Since
$\delta\le\delta_5$, \eqref{eq:scalar} rules out this interval. Thus $x<X$.
Moreover $d+\delta\le1/2$ and $c_1^2/(4h_0)<4Mt$, so the right side of
\eqref{eq:scalar} is at most $m/2+4Mt<(3/5)d_XC_0(1-\delta_5)$.
It follows that $\eta=(x-\delta)/(1-\delta)\le3x/5<1/20$.

Now set $E=0.29$ and
\[
  L=MX\left[4Mt(1-X)+\frac{c_1^2 X}{16Mv(1-X)}\right].
\]
Complete the square in $y$ in \eqref{eq:raw-residual}. Since $x(1-x)$ and
$x^2/(1-x)$ increase on $[0,X]$, we get $\mu^2\|F\|_{-1}^2\le L+MXw$.
Substitution of the rational constants gives
\[
  E^2 d_X^2-L>0.0016,\qquad
  Xm\,\bigl[MX-2E^2 d_X/C_1\bigr]_+<0.0012,
\]
and hence $L+MXw<E^2(d_X+w/C_1)^2\le E^2(d+w/\Cphi)^2$ for $0\le w\le Xm$.
Equation \eqref{eq:W-energy} yields $\norm{\nabla W}<E$. Since $W\perp\phi$
and the second disk eigenvalue is $j_{1,1}^2>14.6$, we have
$\norm{W}\le E/j_{1,1}<0.076$.
\end{proof}

\section{Connectedness at a common interval of heights}\label{sec:levels}
Retain the normalization and spectral bound of Section~\ref{sec:disk}. Let
\[
  r_-=\frac25,\quad r_+=\frac{81}{100},\quad s_\pm=r_\pm^2,\qquad
  E=0.29,\quad Z=0.076,\quad M=5.784,\quad \Lambda=6.03.
\]
Write $\tau=t_*/c$. Orthogonality of $W$ and $\phi$ in energy, the identity
\eqref{eq:V-pde}, and the estimates \eqref{eq:disk-constants} and
\eqref{eq:energy} give
\begin{equation}\label{eq:tau}
  \tau^2=\davg{pV^2}=(1-\delta)\Cphi+\lambda^{-1}\norm{\nabla W}^2,
  \qquad
  0.508<\tau<0.534.
\end{equation}
Also $\kappa\le\mu<M$, $\lambda\le\mu/(1-\delta_5)<\Lambda$,
$2\lambda ae<10/3$ and $\lambda x<1/2$. Set $U=G(\phi h)$.
Expanding the density $p=|a+h|^2$ in the Green representation gives
\begin{equation}\label{eq:green}
  V=(1-\eta)\phi+2a\lambda\Re U+\kappa GW+2a\lambda G(W\Re h)+\lambda G(|h|^2 V).
\end{equation}

\subsection{Inner and outer values and a transverse radial derivative}
Put $B_0(s)=s^2(3-2s)$ and $\Gamma(s)=\pi\norm{G_z}$. Define
\[
  b_2(s)=\frac{1+B_0(s)}{8},\qquad
  b_h(s)=\sqrt{\frac{1+3B_0(s)}{32}},\qquad
  b_\phi(s)=\sqrt{b_2(s)\,G\bigl[B_0(1-3s/8)^2\bigr](z)}.
\]
and
\[
  A_0(s)=\frac{11/40+B_0(s)/5}{\rho},\qquad
  A_+(s)=\frac{25/96+B_0(s)/21}{\rho},\qquad
  T_W(s)=\frac{2ZE}{1+\sqrt\rho}.
\]
Appendix~\ref{app:green} proves the pointwise estimates for $GW$, $U$,
$G(|h|^2)$ and the radial derivatives used below, in particular
\eqref{eq:uniform-G}, \eqref{eq:sources-all} and \eqref{eq:u-slope}.

Take $w(s)=B_0(s)(1-3s/8)^2$. Since $w\ge0$ is radial, $Gw$ is radially
decreasing. Solving for each monomial gives
\[
  Gw(s)=\sum_{j=2}^5[t^j]w(t)\,\frac{1-s^{j+1}}{4(j+1)^2}
       =\frac{1313}{38400}-\frac{s^3}{12}+\frac{17s^4}{256}-\frac{123s^5}{6400}+\frac{s^6}{512}.
\]
Pairing the weighted kernel with $\norm{W}<Z$, or the unweighted kernel
with $\norm{Wh}\le eE/\sqrt\mu$, gives the two bounds
\[
  |G(W\Re h)(z)|\le e\min\left\{Zb_h(s),\ \frac{E\Gamma(s)}{\sqrt\mu}\right\}.
\]
Use $Y(s)=8\Gamma(s)^2$ from \eqref{eq:Y-exact} and $B_0\le9s/8$ to get
\[
  256\bigl(b_h(s)\Gamma(s)\bigr)^2=\bigl(1+3B_0(s)\bigr)Y(s)
    \le\left(1+\frac{27s}{8}\right)(1-s)\le\frac{1225}{864}<\left(\frac65\right)^2.
\]
Let $n_0=\sup_\D(V-\phi)$. It is finite and is nonnegative since $W\perp\phi$.
Subtract $\phi$ from \eqref{eq:green}. Use positivity of the Green kernel,
$b_2\le1/4$, $b_\phi\le\sqrt{Gw(0)/4}$, and the geometric mean of the two
bounds on the mixed term. Writing $n=13/40$, we obtain
\[
  \frac78 n_0\le\frac{10}{3}(0.019)+0.19MZ
    +\frac{10}{3}\sqrt{\frac{3ZE}{40\sqrt{5.783}}}
    +\frac12\sqrt{\frac{1313}{153600}}<\frac78 n,
\]
In the inner disk $r\le r_-$ we drop the nonnegative term
$\lambda G(|h|^2 V)$. Equation \eqref{eq:phi-chord-bounds} and
$\eta<1/20$ give
\[
  V(z)\ge\frac{19}{20}(1-s_-)\left(1-\frac{9s_-}{20}\right)
    -\frac{10}{3}(0.012)-0.19MZ-\frac{10}{3}Zb_h(s_-)>\frac{9}{16}.
\]
In the outer region $r\ge r_+$ we have $b_\phi\le\sqrt{Gw(s_+)/4}<0.0721$
and $\Phi(s)\le p_3(s_+)<0.2545$. Retaining the first three subtracted terms
in \eqref{eq:Y-exact} gives
$8\Gamma^2\le(1-s_+)(1-s_+/2-s_+^2/6-s_+^3/12)$, so $\Gamma<0.158$.
Put $J_+=\Lambda(0.0721+n/4)$. The sum $(1-\eta)\phi+\lambda G(|h|^2V)$
is at most $(1-x)(0.2545)/(1-\delta_5)+xJ_+$, an increasing function of $x$.
At $X=4/49$, using $1/\sqrt\mu<0.416$, this gives
\[
  V(z)\le\frac{10}{3}(0.0152)+\frac{3}{20}ME(1-r_+)
    +\frac{10}{3}E(0.158)(0.416)
    +\frac{1-X}{1-\delta_5}(0.2545)+XJ_+<0.482.
\]

On $r_-\le r\le r_+$ set
$\overline T=2ZE/(1+\sqrt{1-s_+})$ and
$B_*(r)=(5.783/2)r\,b(r^2)$, where $b$ is defined in \eqref{eq:avg-lower}.
Then $-\phi'(r)\ge B_*(r)$. The kernel derivative is locally integrable
and the sources are smooth in the interior. Away from the observation point
the explicit kernel formula controls differentiation under the integral.
Differentiate \eqref{eq:green} radially. Estimate the mixed term by
Cauchy--Schwarz against $|\partial_rG_z|\dd A$. For the last term use
$(\partial_rG_z)_+$ and $0\le V\le\phi+n$. The result is
\begin{equation}\label{eq:slope-unified}
  -\partial_r V(z)\ge F_*(r):=\frac{19}{20}B_*(r)-\frac{10}{3}\cdot\frac{1}{16}
    -\frac{3}{20}ME-\frac{10}{3}\sqrt{\overline T\,A_0(r^2)}
    -\frac12\left(\frac16+nA_+(r^2)\right).
\end{equation}
Since $B_0/\rho=2s^2+s^2/\rho$, both $A$ functions have nonnegative
power-series coefficients in $r^2$. Thus $A_+(r^2)$ is convex in $r$.
The same holds for $\sqrt{A_0(r^2)}$, the $\ell^2$ norm of a sequence
of nonnegative convex monomials. Since
$(rb(r^2))''=r(-87/20+25r^2/8)<0$, the function $F_*$ is concave.
Substitution at $r_-$ and $r_+$ gives lower bounds $1.027$ and $1.385$
for $B_*$. The respective upper bounds for $\sqrt{\overline T A_0}$ are
$0.098$ and $0.185$, and those for $A_+$ are $0.314$ and $0.858$.
These comparisons follow by squaring positive quantities.
Consequently $F_*(r_-)>1/20$ and $F_*(r_+)>3/200$. Throughout the annulus,
\[
                         -\partial_rV>1/100 .
\]

\subsection{Exclusion of a nonconvex minimizer}
Whenever $|t/t_*-1|\le1/20$, \eqref{eq:tau} implies
\[
  0.482<\frac{19}{20}(0.508)<\frac{t}{c}<\frac{21}{20}(0.534)<\frac{9}{16}.
\]
Along every ray $V(re^{i\theta})$ is above $t/c$ for $r\le r_-$, below
$t/c$ for $r\ge r_+$, and strictly decreasing in between. Its unique
crossing has radius $R(\theta)\in(r_-,r_+)$. Interior analyticity and the
implicit function theorem give a real-analytic periodic function $R$.
Consequently $\{V>t/c\}$ is star-shaped about zero. Its conformal image
$\{u>t\}$ is connected.

\begin{proposition}[Convexity of an attained minimizer]\label{prop:attained}
Every global minimizer among simple polygons of prescribed area with at
most $N$ effective sides is convex.
\end{proposition}
\begin{proof}
Suppose an attained minimizer were nonconvex. By Section~\ref{sec:pressure}
it would have $N\ge5$ and $J(P)\le r_5$. Equation \eqref{eq:pent} gives the
spectral bound of Section~\ref{sec:disk}. Choose a regular height in the
intersection of the neighborhood of $t_*$ from Section~\ref{sec:pressure}
and the interval $|t/t_*-1|<1/20$. Equations \eqref{eq:countlow} and
\eqref{eq:countup} give $C(t)\ge2$, while the radial comparison just proved
gives $C(t)=1$.
\end{proof}

\section{Proof of the main theorem and the equality case}\label{sec:assembly}
\begin{proof}[Proof of Theorem~\ref{thm:main}]
Fix $N\ge3$ and $A>0$, and let $m_N$ be the infimum of $J$ for simple
polygons of area $A$ with at most $N$ effective sides. The regular
$N$-gon gives $m_N\le r_N$. If $m_N<r_N$, choose a minimizing sequence
$P_i$ with $J(P_i)<r_N$ for every $i$. By Proposition~\ref{prop:convex}
all its members are nonconvex. A simple polygon has at least three convex
vertices, since its exterior angles sum to $2\pi$ and each positive exterior
angle is strictly less than $\pi$. If some $P_i$ had at least two reflex
vertices, then $N\ge5$ and $p(P_i)\le N-2$. Proposition~\ref{prop:turns}
with $m=N-2$ would give $J(P_i)>r_N$, a contradiction.
Hence each $P_i$ has exactly one reflex vertex. When the side count precludes
even one such vertex, the contradiction is already complete.

Lemma~\ref{lem:compact} gives a simple polygon $Q$ of at most $N$ sides
with $0<|Q|\le A$ and $\lambda_1(Q)\le m_N/A$. Dilate it to a polygon
$Q_A$ of area $A$. By scale invariance,
\[
 m_N\le J(Q_A)=|Q|\lambda_1(Q)\le A\lambda_1(Q)\le m_N .
\]
Thus $Q_A$ attains the global minimum. Proposition~\ref{prop:attained}
gives convexity, and Proposition~\ref{prop:convex} then gives
$m_N=J(Q_A)\ge r_N$, contrary to assumption. This proves \eqref{eq:main}.

Any polygon with $J(P)=r_N$ attains the minimum after rescaling to area
$A$. It is therefore convex, and equality in Proposition~\ref{prop:convex}
forces it to have exactly $N$ effective sides and to be regular.
Conversely the regular $N$-gon attains $r_N$. Its admissibility in the
exactly-$N$ class proves the asserted minimum and uniqueness there as well.
\end{proof}

\appendix
\section{Curvature of the mixed spectral function at acute openings}\label{app:acute}
Here $F$ is the function of Section~\ref{sec:angles}, $s=\alpha/(2\pi)$
parametrizes the opening, and $D=sZ$ denotes an area factor. We use the
sector map and angular decomposition of Wang \cite[Appendix~S1]{wang}, with
the bounds for the first eigenvalue from \cite{unified}.

\subsection{Variation of the mixed eigenvalue}
Put $h=1/5$ and $0<s\le h$. On the unit sector of opening $2\pi s$, a
primitive of $(1-\zeta^{1/s})^{-2s}$ maps the two radii to the equal sides
of a triangle and the circle arc to its base. Put
$\tau=-\log|\zeta|$, $x=\tau/s$, and $\theta=\arg(\zeta)/s$. An overbar
denotes the normalized circle mean in $\theta\in[0,2\pi]$. Integrals with
no measure specified use $\dd\tau\,\dd\theta/(2\pi)$. The positive ground
state is even about the bisector, so its traces on the two sides agree.
Restricting to periodic tests therefore preserves the first eigenvalue.
The trace at $\tau=0$ is zero. Write
\[
 E_s(u)=\int(u_\tau^2+s^{-2}u_\theta^2),\quad
 W=|1-e^{-x+i\theta}|^{-4s},\quad \rho=e^{-2\tau}W,\quad
 L=\inf_{u\ne0}\frac{E_s(u)}{\int\rho u^2}.
\]
Normalize the ground state by $\int\rho u^2=1$. Put $Z=2\int\rho$ and
$D=sZ$. Conformal invariance gives $F(2\pi s)=\pi DL$. The equal sides
have length $s\mathrm B(s,1-2s)$, so their area formula and the gamma
reflection identity give
\[
 A(s):=\log Z=\sum_{\substack{k\ge3\\k\ {\rm odd}}}
       \frac{(2^{k+1}-4)\zeta(k)}{k}s^k .
 \tag{A.1}\label{eq:acute3}
\]
All coefficients are positive. The first two terms of the zeta series with
an integral bound on the tail give $\zeta(3)<5/4$ and $\zeta(k)<17/16$ for
$k\ge5$. Setting $y=2s\le2/5$ gives
\[
 A\le5s^3+\frac{2(17/16)y^5}{5(1-y^2)}<0.05,\qquad
 sA'\le15s^3+\frac{2(17/16)y^5}{1-y^2}<0.15 .
\]
Consequently $Z<1.06$ and $P_s=sD'/D=1+sA'<1.2$. Keeping the cubic
term alone in \eqref{eq:acute3} yields
\[
 \frac{D''}{Ds}=\frac{A''+2A'/s+(A')^2}{s}
 \ge48\zeta(3)>54.
 \tag{A.2}\label{eq:acute4}
\]
The enclosures $3.141<\pi<3.142$ follow from
$\pi/4=4\arctan(1/5)-\arctan(1/239)$ and the alternating arctangent
series. The exponential series with a geometric tail gives $2.718<e<2.719$.

For $j=1,2$ let $W_j$ be the $j$th $s$-derivative at fixed $\tau,\theta$.
Set $W_0=W-1$. Write
\begin{align*}
 m_j&=\int e^{-2\tau}W_j u^2,\qquad
 \beta=s^{-2}\int u_\theta^2/L,\\
 H(\phi)&=E_s(\phi)-L\int\rho\phi^2,\qquad
 f(\phi)=E_s'(u,\phi)-L\int e^{-2\tau}W_1u\phi .
\end{align*}
On the energy-orthogonal complement of $u$, use the bilinear form associated
with $H$ as inner product, and put $G=\|f\|_{H^*}^2$ on that complement.
Differentiation gives $L'/L=-2\beta/s-m_1$. The vector
$u'+m_1u/2$ belongs to the complement and represents $-f$ there.
Differentiating a second time gives
\[
 \frac{(DL)''}{DL}=\frac{D''}{D}-m_2-2P_sm_1/s+2m_1^2+
       \frac{\beta}{s^2}(6-4P_s+4sm_1)-2G/L .
 \tag{A.3}\label{eq:acute5}
\]
Indeed $E_s''(u,u)=6L\beta/s^2$ and $f(u)=L'$. Substitution in the second
variation gives the displayed identity. The differentiations are valid
locally in $s>0$ in the indicated range. Near $\tau=0$ use
$\av{\phi^2}(\tau)\le\tau E_s(\phi)$ and the bound
$C\tau^{-4s}(1+|\log\tau|)^2$ for the weight and its first two
derivatives. At infinity the mass weight decays exponentially. These
bounds give twice continuous differentiability in form norm. Simplicity
and the spectral gap below yield local twice differentiability of the
normalized eigenpair.

\subsection{The ground state, the gap and multiplier bounds}
Write $U=\av u$ and $v=u-U$. The binomial series gives $\av W\ge1$.
The radial disk test $(1-r^2)(1-r^2/3)$ has Rayleigh quotient $110/19$
by direct integration, whence $L<5.8$. The Green kernel of the radial
equation in $\tau$ is $\min(\tau,y)$. Thus Cauchy--Schwarz gives
\[
       0\le U(\tau)\le L\tau\int\rho u
        \le L\sqrt{Z/2}\,\tau<C\tau,\qquad C=17/4 .
 \tag{A.4}\label{eq:acute6}
\]
For a function of circle mean zero, apply the ground-state identity in
every nonzero Fourier mode with the positive comparison function
$x^{4/5}e^{-|k|x}$. The result is
\[
  \int_0^\infty(|\psi_x|^2+k^2|\psi|^2)\dd x
  \ge\int_0^\infty P(x)|\psi|^2\dd x,\qquad P(x)=0.16/x^2+1.6/x .
 \tag{A.5}\label{eq:acute7}
\]
If $x\le1$, then $1-e^{-x}\ge xe^{-x/2}$, so $e^{-2sx}W\le x^{-0.8}$.
If $x\ge1$, use $W\le8/5$ and $xe^{-2sx}\le1/(2es)$. Both ranges give
$Ls^2e^{-2sx}W/P\le.23$. On the angular mean-zero space the mass form
multiplied by $L$ is thus bounded by $\sigma E_s$, $\sigma=0.23$.

Consider the codimension-one space defined by $a(d)=0$, where
$a=\av\phi$ and $d=3/4$. Write $\phi=a+\psi$. On $(0,d)$, the square
of $a(\tau)$ is at most $\tau(1-\tau/d)$ times its energy on this interval,
using its two zero endpoints. On $(d,\infty)$ the corresponding factor
is $\tau-d$. Integrating the unperturbed mass gives an energy coefficient
at most
\[
 L\max\{(7e^{-3/2}-1)/12,\,e^{-3/2}/4\}<0.33 .
\]
Let $a_n(z)=(z)_n/n!$, with $(z)_n$ the rising factorial. Since
$\av W-1=\sum_{n\ge1}a_n(2s)^2 e^{-2nx}$ and $a_n(2s)\le2s$, the
additional radial mass has energy coefficient at most $Ls^4\zeta(2)<0.02$,
using $|a|^2\le\tau E_s(a)$. Similarly,
$\Var_\theta W\le\sum_{n\ge1}a_n(4s)^2 e^{-2nx}\le16s^2/(e^{2x}-1)$.
Equation \eqref{eq:acute7} and $1/P\le5x/8$ bound the mixed coefficient by
\[
 Ls^3\Big(10\int_0^\infty x^2/(e^{2x}-1)\dd x\Big)^{1/2}
 =Ls^3\sqrt{10\zeta(3)/4}<0.09.
\]
The diagonal coefficients of the two-component mass form are less than
$0.35,0.23$ and its off-diagonal coefficient has absolute value less than
$0.09$. Hence its quadratic form is bounded by $E_s/2$ on this
codimension-one space. The min--max principle bounds the second eigenvalue
of the mass form relative to energy by $1/2$. The largest eigenvalue is
one with eigenvector $u$. Consequently
\[
             H\ge E_s/2\quad\text{on }u^{\perp_{E_s}}.
 \tag{A.6}\label{eq:acute8}
\]

For a mean-zero function $v$ and an arbitrary test $\phi$, put
$\|\cdot\|_E=E_s(\cdot)^{1/2}$. We shall use
\begin{align}
 L\int e^{-2\tau}|W_j|v^2&\le T_j s^{2-j}\|v\|_E^2,
   \quad (T_1,T_2)=(12,60),\notag\\
 L\int e^{-2\tau}|W_1v\phi|&\le25s\|v\|_E\|\phi\|_E .
 \tag{A.7}\label{eq:acute9}
\end{align}
To prove these estimates put $b=4s$, $l=-\log(1-e^{-x})$,
$a=x/(e^x-1)$ and $d_0=x^2 e^x/(e^x-1)^2$. Positivity of the
Fourier coefficients gives
\[
 |sW_1|\le g_1=(1-e^{-x})^{-b}b(l+a),\quad
 |s^2W_2|\le g_2=(1-e^{-x})^{-b}\{b^2(l+a)^2+bd_0\}.
\]
By \eqref{eq:acute7} and the one-dimensional Hardy inequality with
weight $1/(4x^2)$ it suffices to bound $Le^{-2sx}g_j/P$ and
$2Le^{-2sx}xg_1/\sqrt P$.

For $x\le1$ write $y=-\log x$. Then $l+a\le3/2+y$, $d_0\le1$ and
$e^{-2sx}(1-e^{-x})^{-b}\le x^{-0.8}$. Apply $1/P\le5x/8$ and
$\sup_{y\ge0} y^m e^{-\eta y}=(m/(e\eta))^m$. The three upper bounds
are less than $10,54,16$, respectively. For example the second is at most
$(5L/8)[0.64(9/4+15/e+100/e^2)+0.8]$. For $x\ge1$, using
$1/(1-e^{-x})<8/5$ gives
\[
 g_1\le2.1(1+x)e^{-x},\qquad
 g_2\le2.7(1+x)^2e^{-2x}+3.3x^2e^{-x}.
\]
The same elementary maximum bound gives upper bounds $8,23$ for the
first two quantities and $2(5.8)(0.8)(2.1)(0.42+0.82)<25$ for the last.
This proves \eqref{eq:acute9}.

\subsection{Circle variances and radial moments}
Define
\[
 N_j(s)^2=s^{-2}\int_0^\infty
       \frac{x^2e^{-4sx}\Var_\theta(s^jW_j)}{P(x)}\,\dd x,\qquad
 e_j=Cs N_j(s),\quad j=0,1,2 .
\]
Changing to $x$ and applying \eqref{eq:acute6}, \eqref{eq:acute7} and
Cauchy--Schwarz shows that $e_js^{5/2-j}$ bounds the energy-dual norm
of $\psi\mapsto\int e^{-2\tau}W_jU\psi$ on the mean-zero space.
We claim
\[
 e_0\le Ch(0.8),\qquad e_1\le Ch(2.7),\qquad s^2e_2\le h^2Ch(13).
 \tag{A.8}\label{eq:acute10}
\]

First replace $1/P$ by $5x/8$. Expand the binomial coefficients as
polynomials in $s$. Every nonzero Fourier mode of $W$ is a sum of positive
terms $c s^a e^{-Nx}$ with $a\ge1$. Indeed, the coefficient of mode $l$
is $\sum_{a_1-a_2=l}a_{a_1}(2s)a_{a_2}(2s)e^{-(a_1+a_2)x}$.
At fixed $\tau$, applying $s^j\partial_s^j$ to one monomial multiplies it
by $a+Nx$ for $j=1$ and by $(a-1+Nx)^2+a-1$ for $j=2$. A product of two
terms of the same mode has even total decay rate $2n$. Let $d$ be its
total degree in $s$ and $g=d/2-1+nx$. The products of the multipliers
are bounded by $M_1=(1+g)^2$ and $M_2=g^2(1+g)^2$, respectively. Indeed,
with $v_i=a_i-1+N_ix\ge0$, we have $(v_1+v_2)/2=g$. The first multiplier
is $1+v_i$ and the second is at most $v_i(1+v_i)$. The two bounds follow
by arithmetic--geometric mean. Set $M_0=1$.

After integration of the positive terms, the bound for $s^2N_j(s)^2$
is a positive linear combination of $s^d/(2n+4s)^{4+m}$, with $d\ge2$
and $m\le2j$. Its individual terms have logarithmic derivative with
respect to $\log s$ at least $2-(4+m)2/7$. The resulting bound increases
in $s$ for $j=0,1$ and increases after multiplication by $s^4$ for $j=2$.
Thus the endpoint $s=h$ suffices.

At $n=1$ the zeroth-order circle variance is exactly $8h^2 e^{-2x}$.
For the differentiated variances its extra factors are $(1+x)^2$ and $x^4$.
With $w=14/5$ the corresponding contributions to $N_j(h)^2$ are at most
\[
 H_0=30/w^4,\quad H_1=5(6/w^4+48/w^5+120/w^6),\quad H_2=25200/w^8 .
\]
For $n\ge2$ we can bound the nonzero modes by all modes. Put $z=s/h$.
The $e^{-2nx}$ coefficient in the circle mean of $W^2$ is
$a_n(0.8z)^2$, and $a_n(0.8)\le a_2(0.8)=0.72$. Normalize the coefficients
of this polynomial as probabilities. The random degree is $d=2+Y$, where
$Y$ is a sum of independent Bernoulli variables, two for each
$1\le i<n$, each with success probability $0.8/(i+0.8)$. Since
$1/(i+0.8)\le1/(i+1/2)\le\int_i^{i+1}\dd y/y$, we have
$\mathbb EY\le1.6\log n$. For $j=0,1,2$ take, respectively,
\[
 (\ell,K,\nu)=(0,1,0),\quad (1/2,18/5,1),\quad (1,8,2).
\]
Then $M_j(g)\le K e^{\ell g}$. For $j=2$ use
$g^2(1+g)^2\le(g+1/2)^4$. Independence gives
$\mathbb E e^{\ell Y/2}\le\exp[(e^{\ell/2}-1)\mathbb EY]\le n^\nu$.
Put $b_*=2-\ell$ and use $\int_0^\infty x^3 e^{-\lambda x}\dd x=6/\lambda^4$.
Retain the terms $n=2,3$ and bound the remaining tail by
$n^\nu/(b_*n+0.8)^4\le n^{\nu-4}/b_*^4$ and the midpoint integral from
$7/2$. This gives
\[
 N_j(h)^2\le H_j+48.6K\left[
 \frac{2^\nu}{(2b_*+0.8)^4}+\frac{3^\nu}{(3b_*+0.8)^4}
 +\frac{(7/2)^{\nu-3}}{(3-\nu)b_*^4}\right].
 \tag{A.9}\label{eq:acute11}
\]
For the three triples the right sides are less than $0.63,6.90,161$.
Taking square roots proves \eqref{eq:acute10}. The exact leading factor
$x^4$ has the same monotonicity after multiplication by $s^4$ when $j=2$.

For the radial moments define
$J_{mj}=s^{-(m+3-j)}\int_0^\infty\tau^m e^{-2\tau}\av{W_j}\dd\tau$,
$j=1,2$. We need
\[
 s^2 J_{11}\le6h^2,\qquad s^2 J_{21}\le4h^2,\qquad
 s^2 J_{22}\le15h^2 . \tag{A.10}\label{eq:acute12}
\]
On expanding the positive series for $\av W$, each term of $s^2J_{mj}$
is proportional to $s^d/(2n+2s)^{m+k+1}$ with $d\ge2$ and $k\le j$.
It increases with $s$ since $s/(n+s)\le1/6$ and $m+j+1\le5$ for the
three moments. At $s=h$ put $D_n=\sum_{i=0}^{n-1}0.4/(i+0.4)$ and
$U_n=\sum_{i=0}^{n-1}(0.4/(i+0.4))^2$. The $n$th contribution to
$J_{mj}(h)$ is
\[
 \frac{a_n(0.4)^2}{h^2}\sum_{k=0}^j \frac{b_{jk}(m+k)!}{[2(n+h)]^{m+k+1}},
 \quad (b_{1k})=(2D_n,2n),\quad
 (b_{2k})=(4D_n^2-2U_n,\,4n(2D_n-1),\,4n^2).
\]
The first terms for $J_{11},J_{21},J_{22}$ are $275/108,125/48,2875/324$.
For $n\ge2$ we have $a_n(0.4)^2/h^2<2$ and $D_n\le\sqrt n$, the latter
by induction on consecutive terms. The three tails are therefore bounded by
\[
 \sum(n^{-3/2}+n^{-2}),\qquad
 \sum(n^{-5/2}+1.5n^{-3}),\qquad
 2\sum(n^{-2}+3n^{-5/2}+3n^{-3})
\]
respectively, with each sum starting at $n=2$.
Now $\sum_{n\ge2}n^{-a}\le(3/2)^{1-a}/(a-1)$ gives \eqref{eq:acute12}.

\subsection{The curvature bound}
Let $V=e_0/(1-\sigma)<0.9$. On mean-zero tests the ground-state equation has
angular source $W_0U+Wv$, with the common radial weight $e^{-2\tau}$.
Equations \eqref{eq:acute7} and \eqref{eq:acute10} and the mass bound give
\[
 \|v\|_E\le LV s^{5/2},\qquad e_1<2.3,\qquad s^2e_2<11.1h^2.
\]
Expand $u^2=U^2+2Uv+v^2$ and apply \eqref{eq:acute6}, \eqref{eq:acute10}
and \eqref{eq:acute9} to the respective terms. The result is
\[
 |m_j|\le s^{5-j}\{ C^2J_{2j}+2LV e_j+LT_jV^2s^2\}.
 \tag{A.11}\label{eq:acute13}
\]
Substituting \eqref{eq:acute12} and \eqref{eq:acute13} with $s\le h$ gives
$s|m_1|<0.04$. In particular $6-4P_s+4sm_1>0$ in \eqref{eq:acute5}.

In estimating $f$, the term $E_s'(u,\cdot)$ acts only on the angular
component and has energy-dual norm at most $2\|v\|_E/s$. In the mass
derivative the angular source from $U$ is bounded by $e_1$. For its radial
source use $|\av\phi|\le\sqrt\tau\,\|\av\phi\|_E$ and Cauchy--Schwarz on
the two adjacent moments. These two parts are energy orthogonal. Apply the
mixed estimate in \eqref{eq:acute9} to $W_1v$. This yields
\[
 \|f\|_{E^*}\le Ls^{3/2}
 \left[\sqrt{(2V+e_1)^2+R^2}+25s^2 V\right],\qquad
 R=Cs^2\sqrt{J_{11}J_{21}}<0.84 .
\]
By \eqref{eq:acute8} and \eqref{eq:acute13}, the total loss from all
potentially negative terms in \eqref{eq:acute5}, divided by $s$, is at most
\begin{align*}
 &h^2\left[C^2(15+2.4\cdot4)+2(5.8)(0.9)(11.1+2.4\cdot2.3)
  +5.8(60+2.4\cdot12)(0.9)^2h^2\right]\\
 &\hspace{10mm}{}+4(5.8)h^2
       \left[\sqrt{(1.8+2.3)^2+0.84^2}+0.9\right]^2 < 26+25=51 .
\end{align*}
Combining this with \eqref{eq:acute4} gives $(DL)''>3sDL>0$ and hence
$F''(\alpha)>0$ throughout $0<\alpha\le2\pi/5$.

\section{Pointwise and radial estimates for Green potentials on the disk}
\label{app:green}
We prove the estimates used in Section~\ref{sec:levels}. The subscript $z$
specifies the observation point. Integrals against $\dd A$ use ordinary
area, while $\norm{\cdot}$ uses area divided by $\pi$.

\subsection{The kernel on the ground-state complement}
For $W\in H_0^1(\D)$ with $W\perp\phi$ we claim
\begin{equation}\label{eq:uniform-G}
  |\partial_r GW(z)|\le\frac{3}{20}\norm{\nabla W},
  \qquad
  |GW(z)|\le\min\left\{0.19\norm{W},\ \frac{3}{20}(1-r)\norm{\nabla W}\right\}.
\end{equation}
Use the disk automorphism $w=(z+\zeta)/(1+\bar z\zeta)$. Then
$G_z(w)=-(2\pi)^{-1}\log|\zeta|$ and its area Jacobian is
$\rho^2/|1+\bar z\zeta|^4$. Angular integration of the power series gives
\begin{equation}\label{eq:Y-exact}
  Y(s):=8\pi^2\norm{G_z}^2
   =\rho^2\sum_{k=0}^\infty\frac{s^k}{k+1}
   =\rho\left(1-\sum_{k=1}^\infty\frac{s^k}{k(k+1)}\right)
   \le\rho\left(1-\frac{s}{2}\right)\le\Phi(s),
\end{equation}
where the last inequality uses \eqref{eq:phi-chord-bounds}. Projecting
$\pi G_z$ onto $\phi^\perp$ in $L^2$ reduces the squared norm from $Y/8$
by $\Phi(s)^2/(\mu^2\Cphi)$. Maximizing $Y/8-Y^2/(\mu^2\Cphi)$ gives the
upper bound $\mu^2\Cphi/256<0.19^2$.

For the radial derivative we may take $z=r>0$. Put
$\mathcal D(s)=\|\pi\partial_rG_r\|_{-1}^2$. Restriction to $\phi^\perp$
in the energy inner product gives the squared functional norm
\begin{equation}\label{eq:proj-D}
  \mathcal D(s)-\frac{\phi'(r)^2}{\mu^3\Cphi}
    =\mathcal D(s)-\frac{sB(s)^2}{4\mu\Cphi},
  \qquad
  B(s)=\int_0^1\Phi(st)\dd t,
\end{equation}
since $-\phi'(r)=\mu rB(s)/2$. We have $B'''\le0$, so Hermite
interpolation at $0,0,1$ gives
\begin{equation}\label{eq:avg-lower}
  B(s)\ge 1-\frac{\mu}{8}s(1-s)+\bigl(B(1)-1\bigr)s^2
      \ge b(s):=1-\frac{29}{40}s+\frac{5}{32}s^2,
\end{equation}
Here $\mu/8<29/40$ and $B(1)=2\sqrt{\Cphi/\mu}>69/160$.
Appendix~\ref{app:series} proves the signed expansion
\begin{equation}\label{eq:D-series}
  32\mathcal D(s)=\frac{7}{12}+\frac{79}{36}s-\sum_{k=2}^\infty d_k s^k,
  \qquad
  d_k=\frac{2\bigl(3k^5+23k^4+53k^3+45k^2+14k+6\bigr)}{k^2(k-1)(k+1)^2(k+2)^2(k+3)}>0.
\end{equation}
Keeping the negative terms only through degree six gives an upper bound.
In view of $8/(\mu\Cphi)>128/25$, it suffices to prove positivity on
$[0,1]$ of the polynomial
\[
  Q(s)=32\left(\frac{3}{20}\right)^2-\frac{7}{12}-\frac{79}{36}s
    +\sum_{k=2}^6 d_k s^k+\frac{128}{25}s\,b(s)^2.
\]
For a polynomial of degree at most $m$ the Bernstein representation is
\begin{equation}\label{eq:bernstein}
  q(s)=\sum_{i=0}^m\beta_i\binom{m}{i}s^i(1-s)^{m-i},
  \qquad
  \beta_i=\sum_{j=0}^i\frac{\binom{i}{j}}{\binom{m}{j}}[s^j]q(s).
\end{equation}
At $m=6$ the coefficients of $Q$ satisfy
$1000(\beta_0,\ldots,\beta_6)>(136,624,667,492,251,55,24)$ componentwise
by substitution. This proves the derivative bound in \eqref{eq:uniform-G}.
The disk embedding $H_0^1\hookrightarrow L^p$ holds for every finite $p$,
and elliptic regularity gives $GW\in W^{2,p}\hookrightarrow C^1$ for $p>2$.
Thus point evaluation extends by density. Integrating the derivative
radially from the boundary gives the bound containing $1-r$.

\subsection{Quadratic and mixed sources}
Retain the functions $B_0,b_2,b_h,b_\phi,A_0,A_+,T_W$ from
Section~\ref{sec:levels}. Write $K=|\partial_rG_z(w)|$ and
$K_+=(\partial_rG_z(w))_+$. We prove the estimates
\[
 G(|h|^2)\le x b_2,\qquad \pi\norm{G_z h}\le e b_h,\qquad
 G(\phi|h|^2)\le x b_\phi ,
\]
together with
\begin{equation}\label{eq:sources-all}\begin{aligned}
  \int_\D KW^2\dd A &\le T_W(s), &
  \int_\D K|h|^2\dd A &\le xA_0(s),\\
  \int_\D K_+|h|^2\dd A &\le xA_+(s), &
  \int_\D K_+\phi|h|^2\dd A &\le\frac{x}{6}.
\end{aligned}\end{equation}
Cauchy--Schwarz on the power series for $h$ yields
$\rho^2|h(z)|^2\le x B_0(s)$. Under the automorphism above,
$H=(h\circ w)w'$ has $\norm H=e$ and $|H(0)|^2\le xB_0$.
In the normalized Bergman basis $\sqrt{k+1}\,\zeta^k$ the diagonal
weights of $G(|H|^2)(0)$ and $\pi^2\norm{G_0H}^2$ are
$1/[4(k+1)]$ and $1/[8(k+1)^2]$, respectively. Separating the constant
mode gives the first two estimates. For the third use Cauchy--Schwarz
with measure $G_z\dd A$ and
$\phi^2|h|^2\le xB_0(1-3s/8)^2$.

The kernel formula gives
$|\nabla_zG_z(w)|=(1-|w|^2)/(2\pi|w-z|\,|1-\bar zw|)$.
For $t=|w|$, we have
\[
 2(1-rt)-(1+\sqrt\rho)(1-t^2)
 =[(1+\sqrt\rho)t-r]^2/(1+\sqrt\rho)\ge0
\]
and $\operatorname{div}_w((w-z)/|w-z|)=1/|w-z|$.
Integration by parts against $W^2$ followed by Cauchy--Schwarz yields
$\int_\D KW^2\dd A\le2\norm W\norm{\nabla W}/(1+\sqrt\rho)$.

Assume $x>0$. Rotate to $z=r$ and write $\zeta=t e^{i\theta}$.
Then $2\pi\partial_rG_r(w)=\rho^{-1}(1/t-t)\cos\theta$, with the derivative
taken before the change of variables. For $c(\theta)=|\cos\theta|$ or
$(\cos\theta)_+$ let $K_c=K$ or $K_+$, respectively. In the normalized
Bergman coefficients of $H$, the real symmetric matrix of
$\rho\int_\D K_c|h|^2\dd A$ is
\[
  D_{ij}=\widehat c(i-j)\,\frac{2\sqrt{(i+1)(j+1)}}{(i+j+1)(i+j+3)}
  \qquad(i,j\ge0),
\]
where $\widehat c(m)=(2\pi)^{-1}\int_{-\pi}^{\pi}c(\theta)e^{-im\theta}\dd\theta$.
The entries in the following table are upper bounds.
\[
  \begin{tabular}{cccc}
    \toprule
    $c$ & $|D_{00}|$ & $\|(D_{0j})_{j\ge1}\|_{\ell^2}$ & $\|(D_{ij})_{i,j\ge1}\|_{\ell^2\to\ell^2}$ \\
    \midrule
    $|\cos\theta|$ & $17/40$ & $1/20$ & $9/40$ \\[2pt]
    $(\cos\theta)_+$ & $3/14$ & $3/32$ & $1/6$ \\
    \bottomrule
  \end{tabular}
\]
For verification, set $R_j=(j+2)/[(j+1)(j+3)]$. The radial factor
in $D_{ij}$ is at most $R_{i+j}$ by arithmetic--geometric mean, and
$R_j$ decreases in $j$. For $(\cos\theta)_+$ the absolute Fourier
coefficients at $0,\pm1,\pm2l$ are $1/\pi,1/4,1/[\pi(4l^2-1)]$ for
$l\ge1$, with the remaining odd coefficients zero. For $|\cos\theta|$
the even coefficients double and the odd ones vanish. Since
$\sum_{l\ge1}(4l^2-1)^{-1}=1/2$, absolute row sums of the block
$i,j\ge1$ for $|\cos\theta|$ are bounded by
$\max\{(2R_2+R_4)/\pi,4R_4/\pi\}<9/40$, treating $i\le2$ and $i\ge3$
separately. For $(\cos\theta)_+$ the bounds at $i=1$, at $i=2$, and for
$i\ge3$ are respectively
\[
  \frac{R_2+R_4/2}{\pi}+\frac{R_3}{4},\qquad
  \frac{R_4+R_6/2}{\pi}+\frac{R_3+R_5}{4},\qquad
  \frac{R_6+R_4}{\pi}+\frac{R_5+R_7}{4},
\]
All three are less than $1/6$, and Schur's test gives the operator
bounds. In row zero the radial factor is $2/[(j+3)\sqrt{j+1}]$.
Using $\sum_{l\ge2}(4l^2-1)^{-2}\le1/90$, the squared norm of this row
after deleting its first entry is at most
$4[4/(75\cdot9)+4/(245\cdot90)]/\pi^2<(1/20)^2$ for $|\cos\theta|$.
For the positive part the corresponding bound is $1/128+1/40^2<(3/32)^2$.

Put $m_0=|H(0)|^2/x\le B_0$. If the table entries are $d_0,d_1,d_2$,
the normalized quadratic form is at most
\[
 d_0m_0+2d_1\sqrt{m_0(1-m_0)}+d_2(1-m_0)\le d_2+d_1+(d_0-d_2)m_0 .
\]
Substitution gives the bounds involving $A_0,A_+$. Where $K_+>0$ one has
$\Re\zeta>0$, $|w|>r$ and $1-|w|^2\le\rho(1-t^2)$. Hence
$\phi(w)\le\rho(1-t^2)(1-3s/8)$. The extra radial weight multiplies the
matrix majorant entry by $4/(i+j+5)$, at most $4/5,2/3,4/7$ on the
three respective blocks. Thus
\[
  \frac{1}{x}\int_\D K_+\phi|h|^2\dd A
    \le\left(1-\frac{3s}{8}\right)\left(\frac{53}{336}+\frac{8B_0(s)}{105}\right)
    \le\left(1-\frac{3s}{8}\right)\left(\frac{53}{336}+\frac{3s}{35}\right)<\frac16,
\]
The last step uses $B_0\le9s/8$ and the maximum of a quadratic. This
proves \eqref{eq:sources-all}. The case $x=0$ follows from $h=0$.

\subsection{Cancellation in the linear potential}
For $U=G(\phi h)$ we have $|U(z)|\le e\,v(s)$, where
\begin{equation}\label{eq:v-linear}
  v(s)=\sqrt3\,s\sqrt{(1-s)(1-s/4)}\,c_0(s),
  \qquad
  c_0(s)=\frac{354-286s+59s^2-5s^3}{7680}.
\end{equation}
We shall prove $v<0.019$ throughout $[0,1]$, $v<0.012$ on $s\le s_-$ and
$v<0.0152$ on $s\ge s_+$. On $[s_-,s_+]$ we also have
\begin{equation}\label{eq:u-slope}
  |\partial_r U(z)|\le\frac{e}{16}.
\end{equation}

Since $(k+1)u^k\dd u$ moves to the right as $k$ increases and $\Phi$ is
decreasing, \eqref{eq:free-profile} gives
$(k+1)P_{k,\Phi}\le3P_{2,p_3}=3\rho c_0$ for $k\ge2$.
Moreover $\sum_{k\ge2}s^k/(k+1)\le s^2(1-s/4)/(3\rho)$.
Cauchy--Schwarz on the series for $U$ gives the stated $e\,v(s)$ bound.
The cubic $c_0$ is decreasing and convex, with $c_0(1)\le3c_0(0)/8$.
Use its endpoint chord and arithmetic--geometric mean to get
\[
  v(s)\le\sqrt3\,c_0(0)\,s\left(1-\frac{5s}{8}\right)^2
      \le\frac{32\sqrt3}{135}\,c_0(0)<0.019.
\]
In the inner disk use $s\sqrt{1-s}\le s_-\sqrt{1-s_-}$ and bound the two
remaining factors by their values at zero. In the outer region use
$s\sqrt{1-s}\le2/(3\sqrt3)$ and bound the remaining decreasing factors
at $s_+$. This gives the sharper constants on the two regions.

To preserve the cancellation in the derivative of each mode, put
\[
  A(s)=\int_s^1\Phi(t)\dd t,\qquad
  B_k(s)=\int_0^1 u^k\bigl[\Phi(u)+s\Phi(su)\bigr]\dd u.
\]
Summing \eqref{eq:free-profile} gives
$4U(z)=\int_s^1\int_0^1\Phi(tu)h(uz)\dd u\dd t$. Differentiate, replace
$r\partial_rh(uz)$ by $u\partial_uh(uz)$ and integrate by parts in $u$.
Using $\partial_t[t\Phi(tu)]=\Phi(tu)+tu\Phi'(tu)$ yields
\[
  4r\partial_r U(z)=A(s)h(z)-\int_0^1\bigl[\Phi(u)+s\Phi(su)\bigr]h(uz)\dd u.
\]
Thus for $e>0$, $16|\partial_r U|^2/e^2\le
\sum_{k\ge2}(k+1)s^{k-1}(A-B_k)^2$.
Integrating \eqref{eq:phi-chord-bounds} for $s_-\le s\le s_+<2/3$ gives
\[
  \frac{13}{40}\rho^2\le\rho^2\left(\frac{11}{40}+\frac{3\rho}{20}\right)
    \le A(s)\le\rho^2\left(\frac{5}{16}+\frac{5\rho}{48}+\frac{\rho^2}{64}\right)
    \le\frac{5}{12}\rho^2.
\]
Using $\Phi(u)\le(1-u)(1-3u/8)$ also yields
\[
  B_k(s)\le c_k+\frac{s}{k+1}-\left(\frac{11}{8(k+2)}-\frac{1}{4(k+3)}\right)s^2,
  \qquad
  c_k=\frac{1}{k+1}-\frac{11}{8(k+2)}+\frac{3}{8(k+3)}.
\]
Completion of the square in $s$ gives upper bounds
$10771/67680,117/1120,41549/546000$ for $B_2,B_3,B_4$, respectively.
In particular $B_2\le b_2'=4/25$, $B_3\le b_3'=21/200$, and for $k\ge4$,
$B_k\le B_4<(13/20)(1-s_+)^2\le2A$. Put $g=(13/20)\rho^2\le2A$.
For $k=2,3$ we have $(A-B_k)^2\le A^2+b_k'(b_k'-g)_+$, with the
bound $A^2$ for the remaining modes. Set
\[
  P_0(s)=\frac{25}{144}s(3-2s)\rho^2,\quad
  P_1(s)=P_0(s)+3sb_2'(b_2'-g),\quad
  P_2(s)=P_1(s)+4s^2 b_3'(b_3'-g).
\]
Since $b_2'>b_3'$, summing gives
$16|\partial_rU|^2/e^2\le\max(P_0,P_1,P_2)$. We used
$\sum_{k\ge2}(k+1)s^{k-1}=s(3-2s)/\rho^2$. The inequality
$(3-2s)\rho^2\le3(1-8s/9)^3$ gives
$P_0\le25\cdot243/(48\cdot2048)<1/16$. For $i=1,2$ use
\eqref{eq:bernstein} to expand $1/16-P_i(2u/3)$ in the degree-four
Bernstein basis on $[0,1]$. Its coefficients multiplied by $1000$ strictly
exceed $(62,14,23,22,12)$ and $(62,14,40,33,6)$ componentwise,
respectively. Hence $\max P_i<1/16$, proving \eqref{eq:u-slope}.

\section{The signed series for the radial Green energy}\label{app:series}
It remains to prove \eqref{eq:D-series}. Set $s=r^2$, $w=t e^{i\theta}$ and
$F_r(w)=\pi\partial_r G_r(w)$. In the Fourier cosine expansion of the
logarithmic kernel derivative the constant mode is
$-(2r)^{-1}\mathbf 1_{\{t<r\}}$. For $k\ge1$ the coefficient is
\[
  f_k(t)=\frac{1}{2r}\begin{cases}
    -\bigl[(t/r)^k+(tr)^k\bigr], & t<r,\\[3pt]
    (r/t)^k-(rt)^k, & t>r.
  \end{cases}
\]
Write the corresponding mode of the potential as $t^k p_k(t)\cos(k\theta)$.
Then $-(t^{2k+1}p_k')'=t^{k+1}f_k(t)$ with $p_k(1)=0$. Integration by
parts gives its contribution to $\mathcal D(s)$ as
\[
  E_k=\int_0^1 t^{-2k-1}\left(\int_0^t f_k(v)v^{k+1}\dd v\right)^2\dd t,
\]
The constant mode has twice the angular normalization and contributes
$E_0=s(1-2\log s)/32$. Evaluate the inner integral separately for $t<r$
and $t>r$, then use $\xi=s/t^2$ in the second range. This gives
\begin{equation}\label{eq:modal-kernel}
  \frac{32E_k}{s}=\frac{(1+s^k)^2}{(k+1)^2(k+2)}
    +\int_s^1\xi^{k-3}\left(1-\frac{k+2}{k+1}\xi-\frac{(s/\xi)^k}{k+1}\right)^2\dd\xi.
\end{equation}
By angular orthogonality, $\mathcal D(s)=\sum_{k\ge0}E_k$. Write
$L=\sum_{j\ge1}s^j/j$ and $T=\sum_{j\ge1}s^j/j^2$.
For $k\ge3$ integration of \eqref{eq:modal-kernel} gives
$32E_k/s=c_k+s^kh_k(s)$, where
\begin{align*}
  c_k &= \frac{1}{k+2}-\frac{3}{k+1}+\frac{4}{k}-\frac{3}{k-1}+\frac{1}{k-2},\\
  h_k(s) &= \frac{3}{s(k-1)}+\frac{2s^2+1}{s^2(k+2)}
    +\frac{s-2}{s^2(k+1)}+\frac{(s+1)^2}{s^2(k+1)^2}
    -\frac{1}{s^2(k-2)}-\frac{4}{k}.
\end{align*}
The constant terms telescope to $\sum_{k\ge3}c_k=5/12$. Direct integration
of the first two nonconstant modes gives
\[
  \frac{32E_1}{s}=3\log s-\frac{37s}{12}+\frac{17}{6}+\frac{7}{12s},
  \qquad
  \frac{32E_2}{s}=-\log s-\frac{25s^2}{18}+\frac{32s}{9}-\frac{37}{18},
\]
Their logarithmic terms cancel the logarithm in $32E_0/s$. All series
converge absolutely for $0<s<1$, so shifting indices in the sum of
$s^kh_k(s)$ gives
\[
  32\mathcal D(s)=\left(1+\frac{2}{s}+\frac{1}{s^2}\right)T
    +\left(-2s+\frac{3}{s}-\frac{2}{s^2}+\frac{1}{s^3}\right)L
    -\frac{s}{2}-4+\frac{1}{2s}-\frac{1}{s^2}.
\]
The coefficients of $s^{-2},s^{-1}$ vanish. The constant and linear
coefficients are $7/12$ and $79/36$. For $n\ge2$ the negative of the
coefficient of $s^n$ is
\[
  \frac{2}{n-1}-\frac{3}{n+1}+\frac{2}{n+2}-\frac{1}{n+3}
    -\frac{1}{n^2}-\frac{2}{(n+1)^2}-\frac{1}{(n+2)^2}.
\]
Putting these fractions over a common denominator gives the positive
coefficient $d_n$ in \eqref{eq:D-series}. In particular the truncated
series used in Appendix~\ref{app:green} is an upper bound throughout its
stated interval.

\enlargethispage{1.5\baselineskip}
\begin{thebibliography}{9}

\bibitem{brasco}
L. Brasco, \emph{On torsional rigidity and principal frequencies: an
invitation to the Kohler--Jobin rearrangement technique},
ESAIM Control Optim. Calc. Var. \textbf{20} (2014), 315--338.
\href{https://doi.org/10.1051/cocv/2013065}{doi:10.1051/cocv/2013065}.

\bibitem{cheng}
\emph{Convexity of polygonal Dirichlet eigenvalue minimizers}, manuscript
made available by Zhuo Cheng, version \texttt{4a79a2a}, Theorem~1.1,
Sections~2--5 and Appendices~A and~B.
\href{https://github.com/ChengZhuoCZ/polygonal-faber-krahn-convexity/blob/4a79a2ae9a56d9274503b1f099634a6db8e19ffe/convexity_polygonal_minimizers_global.tex}{Archived source at the cited commit}.

\bibitem{SC}
T. A. Driscoll and L. N. Trefethen, \emph{Schwarz--Christoffel Mapping},
Cambridge Monographs on Applied and Computational Mathematics, vol.~8,
Cambridge University Press, Cambridge, 2002.

\bibitem{kato}
T. Kato, \emph{Perturbation Theory for Linear Operators}, Classics in
Mathematics, Springer-Verlag, Berlin, 1995.

\bibitem{ps}
G. P\'olya and G. Szeg\H{o}, \emph{Isoperimetric Inequalities in Mathematical
Physics}, Annals of Mathematics Studies, vol.~27, Princeton University
Press, Princeton, NJ, 1951.

\bibitem{unified}
\emph{A unified analytic proof of the polygonal P\'olya--Szeg\H{o}
conjecture}, working manuscript in Chinese, November 22, 2026, Sections~2--4
and Appendix~A.

\bibitem{wang}
H. Wang, \emph{The P\'olya--Szeg\H{o} inequalities for polygons}, preprint,
2026. \href{https://arxiv.org/abs/2611.22305v1}{arXiv:2611.22305v1}.

\end{thebibliography}
\end{document}
